\documentclass[8pt]{beamer}
\usetheme{Madrid}
\usecolortheme{default}
\usepackage{amsmath, amssymb, graphicx, array, multirow, booktabs, xcolor, tikz}
\usetikzlibrary{shapes,arrows,positioning,fit,backgrounds}

% Compact font settings
\setbeamerfont{title}{size=\Large}
\setbeamerfont{subtitle}{size=\small}
\setbeamerfont{author}{size=\scriptsize}
\setbeamerfont{date}{size=\scriptsize}
\setbeamerfont{frametitle}{size=\large}
\setbeamerfont{framesubtitle}{size=\normalsize}
\setbeamerfont{enumerate item}{size=\footnotesize}
\setbeamerfont{itemize item}{size=\footnotesize}
\setbeamerfont{block title}{size=\footnotesize}

% Compact spacing
\setlength{\itemsep}{0.08ex}
\setlength{\parskip}{0.08ex}
\setlength{\leftmargini}{0.3cm}
\setbeamersize{text margin left=3mm, text margin right=3mm}
\addtobeamertemplate{frame begin}{\vspace{-0.4cm}}{}

% Colors
\definecolor{myblue}{RGB}{0,0,255}
\definecolor{myred}{RGB}{255,0,0}
\definecolor{mygreen}{RGB}{0,128,0}
\definecolor{myorange}{RGB}{255,165,0}
\definecolor{mypurple}{RGB}{128,0,128}
\definecolor{mygray}{RGB}{128,128,128}
\definecolor{mygold}{RGB}{255,215,0}

\title{Natural Language Processing}
\subtitle{100 Questions: Pre-processing, BoW, TF-IDF, Naïve Bayes, Evaluation}
\author{GARP RAI Exam Preparation}
\date{}

\begin{document}
	
	\begin{frame}
		\titlepage
	\end{frame}
	
	% ============================================================
	% SECTION 1: INTRODUCTION TO NLP - QUESTIONS 1-10
	% ============================================================
	
	\begin{frame}{Q1: NLP Definition - Tutorial}
		\fontsize{6.5}{7.5}\selectfont
		\textbf{QUESTION: What is Natural Language Processing (NLP)?}
		
		\vspace{0.1cm}
		\textbf{OPTIONS:}
		\begin{enumerate}
			\item A technique for processing numerical data only
			\item An aspect of machine learning concerned with understanding and analyzing human language, both written and spoken
			\item A method for image recognition
			\item A database management technique
		\end{enumerate}
		
		\vspace{0.1cm}
		\textcolor{myblue}{\textbf{ANSWER: B}}
		
		\vspace{0.1cm}
		\textbf{DETAILED TUTORIAL:}
		\begin{itemize}
			\item \textbf{Definition:} NLP is the intersection of computer science, AI, and linguistics
			\item \textbf{Core purpose:} Enable computers to understand, interpret, and generate human language
			\item \textbf{Key challenge:} Most information exists in unstructured, free text format
			\item \textbf{Alternative names:} Content analysis, text mining, computational linguistics
		\end{itemize}
		
		\textcolor{myred}{\textbf{EXAM TRAP:}} NLP is about \textcolor{myblue}{understanding human language} - not just processing text!
	\end{frame}
	
	\begin{frame}{Q2: NLP vs Numerical Data - Tutorial}
		\fontsize{6.5}{7.5}\selectfont
		\textbf{QUESTION: Why is NLP more challenging than modeling purely numerical data?}
		
		\vspace{0.1cm}
		\textbf{OPTIONS:}
		\begin{enumerate}
			\item Numerical data is always larger
			\item Textual data is often noisy with errors and ambiguities
			\item NLP requires more computing power
			\item Numerical data is more interesting
		\end{enumerate}
		
		\vspace{0.1cm}
		\textcolor{myblue}{\textbf{ANSWER: B}}
		
		\vspace{0.1cm}
		\textbf{DETAILED TUTORIAL:}
		\begin{itemize}
			\item \textbf{Key challenges in NLP:}
			\begin{enumerate}
				\item \textbf{Ambiguity:} Same words have different meanings (polysemy)
				\item \textbf{Context dependence:} Meaning depends on surrounding text
				\item \textbf{Noise:} Spelling errors, abbreviations, colloquialisms
				\item \textbf{Informal language:} Social media, text messages
				\item \textbf{Sarcasm and tone:} Hard for machines to detect
			\end{enumerate}
			\item \textbf{Examples:} "u r gr8" vs "you are great", "lol" in different contexts
		\end{itemize}
		
		\textcolor{myred}{\textbf{EXAM TRAP:}} NLP challenges include \textcolor{myblue}{ambiguity, noise, and context} - not just technical issues!
	\end{frame}
	
	\begin{frame}{Q3: NLP Applications in Finance - Tutorial}
		\fontsize{6.5}{7.5}\selectfont
		\textbf{QUESTION: Which of the following is an NLP application in finance?}
		
		\vspace{0.1cm}
		\textbf{OPTIONS:}
		\begin{enumerate}
			\item Fraud detection using SEC filings
			\item Customer service call routing
			\item Sentiment analysis from social media
			\item All of the above
		\end{enumerate}
		
		\vspace{0.1cm}
		\textcolor{myblue}{\textbf{ANSWER: D}}
		
		\vspace{0.1cm}
		\textbf{DETAILED TUTORIAL:}
		\begin{itemize}
			\item \textbf{Key NLP applications in finance:}
			\begin{enumerate}
				\item \textbf{Fraud detection:} SEC used NLP to detect accounting fraud
				\item \textbf{Call routing:} Keyword recognition for customer inquiries
				\item \textbf{Sentiment analysis:} Social media monitoring
				\item \textbf{Readability assessment:} Fog index measurement
				\item \textbf{Document similarity:} Assessing new information content
			\end{enumerate}
		\end{itemize}
		
		\textcolor{myred}{\textbf{EXAM TRAP:}} NLP has \textcolor{myblue}{many financial applications} - fraud detection, sentiment, customer service!
	\end{frame}
	
	\begin{frame}{Q4: NLP vs Manual Reading - Tutorial}
		\fontsize{6.5}{7.5}\selectfont
		\textbf{QUESTION: What are the main benefits of NLP over manual document reading?}
		
		\vspace{0.1cm}
		\textbf{OPTIONS:}
		\begin{enumerate}
			\item NLP is always more accurate than humans
			\item Superior speed, no scope to miss aspects, and consistent application
			\item NLP is cheaper than human labor
			\item NLP doesn't make mistakes
		\end{enumerate}
		
		\vspace{0.1cm}
		\textcolor{myblue}{\textbf{ANSWER: B}}
		
		\vspace{0.1cm}
		\textbf{DETAILED TUTORIAL:}
		\begin{itemize}
			\item \textbf{Key advantages of NLP:}
			\begin{enumerate}
				\item \textbf{Speed:} Processes thousands of documents quickly
				\item \textbf{Completeness:} No scope to miss any aspects
				\item \textbf{Consistency:} All documents considered identically
			\end{enumerate}
		\end{itemize}
		
		\textcolor{myred}{\textbf{EXAM TRAP:}} NLP advantages = \textcolor{myblue}{speed, completeness, and consistency}!
	\end{frame}
	
	\begin{frame}{Q5: NLP Data Sources - Tutorial}
		\fontsize{6.5}{7.5}\selectfont
		\textbf{QUESTION: What are the most common NLP data sources in finance?}
		
		\vspace{0.1cm}
		\textbf{OPTIONS:}
		\begin{enumerate}
			\item Only financial statements
			\item Corporate disclosures, social media posts, and news wires
			\item Only social media
			\item Only news articles
		\end{enumerate}
		
		\vspace{0.1cm}
		\textcolor{myblue}{\textbf{ANSWER: B}}
		
		\vspace{0.1cm}
		\textbf{DETAILED TUTORIAL:}
		\begin{itemize}
			\item \textbf{Primary data sources:}
			\begin{enumerate}
				\item \textbf{Corporate disclosures:} 10-Q, 10-K, prospectuses
				\item \textbf{Social media:} X (Twitter), Reddit, Stocktwits, Seeking Alpha
				\item \textbf{News wires:} Wall Street Journal, Dow Jones newswire
			\end{enumerate}
			\item \textbf{Data collection:} API, web scraping with Python
			\item \textbf{Corpus:} Collection of all retrieved documents
		\end{itemize}
		
		\textcolor{myred}{\textbf{EXAM TRAP:}} Common sources = \textcolor{myblue}{corporate disclosures, social media, and news wires}!
	\end{frame}
	
	\begin{frame}{Q6: JPMorgan COiN Platform - Tutorial}
		\fontsize{6.5}{7.5}\selectfont
		\textbf{QUESTION: What did JPMorgan Chase's COiN platform achieve?}
		
		\vspace{0.1cm}
		\textbf{OPTIONS:}
		\begin{enumerate}
			\item It made no significant impact
			\item Significantly reduced time required to review legal documents
			\item It replaced all human reviewers
			\item It only worked for English documents
		\end{enumerate}
		
		\vspace{0.1cm}
		\textcolor{myblue}{\textbf{ANSWER: B}}
		
		\vspace{0.1cm}
		\textbf{DETAILED TUTORIAL:}
		\begin{itemize}
			\item \textbf{COiN (Contract Intelligence) Platform:}
			\begin{itemize}
				\item Automated review of legal documents
				\item Focus on commercial loan agreements
			\end{itemize}
			\item \textbf{Outcomes:}
			\begin{itemize}
				\item Significantly reduced review time
				\item Improved consistency in information extraction
				\item Improved accuracy in information extraction
			\end{itemize}
		\end{itemize}
		
		\textcolor{myred}{\textbf{EXAM TRAP:}} COiN reduced \textcolor{myblue}{legal document review time} significantly!
	\end{frame}
	
	\begin{frame}{Q7: Other NLP Use Cases - Tutorial}
		\fontsize{6.5}{7.5}\selectfont
		\textbf{QUESTION: Which institutions use NLP for specific applications?}
		
		\vspace{0.1cm}
		\textbf{OPTIONS:}
		\begin{enumerate}
			\item Only JPMorgan Chase uses NLP
			\item HSBC (fraud detection), BoA (chatbots), BlackRock (sentiment analysis)
			\item No financial institutions use NLP
			\item Only the SEC uses NLP
		\end{enumerate}
		
		\vspace{0.1cm}
		\textcolor{myblue}{\textbf{ANSWER: B}}
		
		\vspace{0.1cm}
		\textbf{DETAILED TUTORIAL:}
		\begin{itemize}
			\item \textbf{NLP use cases in banking:}
			\begin{enumerate}
				\item \textbf{Fraud detection:} HSBC
				\item \textbf{Customer service:} Bank of America (BoA) chatbots
				\item \textbf{Market sentiment:} BlackRock
			\end{enumerate}
			\item \textbf{Early example:} SEC used NLP to detect accounting fraud
		\end{itemize}
		
		\textcolor{myred}{\textbf{EXAM TRAP:}} HSBC = \textcolor{myblue}{fraud detection}, BoA = \textcolor{myblue}{chatbots}, BlackRock = \textcolor{myblue}{sentiment analysis}!
	\end{frame}
	
	\begin{frame}{Q8: Corpus Definition - Tutorial}
		\fontsize{6.5}{7.5}\selectfont
		\textbf{QUESTION: What is a corpus in the context of NLP?}
		
		\vspace{0.1cm}
		\textbf{OPTIONS:}
		\begin{enumerate}
			\item A single document
			\item The collection of all documents retrieved for analysis
			\item A type of machine learning algorithm
			\item A data pre-processing technique
		\end{enumerate}
		
		\vspace{0.1cm}
		\textcolor{myblue}{\textbf{ANSWER: B}}
		
		\vspace{0.1cm}
		\textbf{DETAILED TUTORIAL:}
		\begin{itemize}
			\item \textbf{Definition:}
			\begin{itemize}
				\item From Latin "corpus" meaning "body"
				\item Collection of all documents available for analysis
				\item May include hundreds or thousands of documents
			\end{itemize}
			\item \textbf{Characteristics:}
			\begin{itemize}
				\item All documents in the corpus are processed together
				\item Vocabulary built from the corpus
			\end{itemize}
		\end{itemize}
		
		\textcolor{myred}{\textbf{EXAM TRAP:}} Corpus = \textcolor{myblue}{collection of all documents} - not a single document!
	\end{frame}
	
	\begin{frame}{Q9: NLP Pre-processing Importance - Tutorial}
		\fontsize{6.5}{7.5}\selectfont
		\textbf{QUESTION: Why is data pre-processing important in NLP?}
		
		\vspace{0.1cm}
		\textbf{OPTIONS:}
		\begin{enumerate}
			\item It's not important
			\item To convert raw text into a format suitable for analysis
			\item To make documents shorter
			\item To remove all words
		\end{enumerate}
		
		\vspace{0.1cm}
		\textcolor{myblue}{\textbf{ANSWER: B}}
		
		\vspace{0.1cm}
		\textbf{DETAILED TUTORIAL:}
		\begin{itemize}
			\item \textbf{Purpose of pre-processing:}
			\begin{enumerate}
				\item Convert text to machine-readable format
				\item Remove non-textual information (HTML, tags)
				\item Standardize the text
				\item Reduce noise and improve accuracy
			\end{enumerate}
			\item \textbf{Initial cleaning:} Extract text from PDFs, remove HTML, handle emojis
		\end{itemize}
		
		\textcolor{myred}{\textbf{EXAM TRAP:}} Pre-processing converts \textcolor{myblue}{raw text to analyzable format}!
	\end{frame}
	
	\begin{frame}{Q10: Document Similarity - Tutorial}
		\fontsize{6.5}{7.5}\selectfont
		\textbf{QUESTION: What is a practical use of document similarity in finance?}
		
		\vspace{0.1cm}
		\textbf{OPTIONS:}
		\begin{enumerate}
			\item To make documents longer
			\item To assess whether a news article constitutes new information
			\item To delete duplicate documents
			\item To remove all documents
		\end{enumerate}
		
		\vspace{0.1cm}
		\textcolor{myblue}{\textbf{ANSWER: B}}
		
		\vspace{0.1cm}
		\textbf{DETAILED TUTORIAL:}
		\begin{itemize}
			\item \textbf{Document similarity applications:}
			\begin{enumerate}
				\item \textbf{New information assessment:} Determine if news article provides new information
				\item \textbf{Disclosure tracking:} Assess if tone of accounting disclosures has changed
			\end{enumerate}
		\end{itemize}
		
		\textcolor{myred}{\textbf{EXAM TRAP:}} Document similarity helps assess \textcolor{myblue}{new information content}!
	\end{frame}
	
	% ============================================================
	% SECTION 2: DATA PRE-PROCESSING - QUESTIONS 11-25
	% ============================================================
	
	\begin{frame}{Q11: Tokenization Definition - Tutorial}
		\fontsize{6.5}{7.5}\selectfont
		\textbf{QUESTION: What is tokenization in NLP pre-processing?}
		
		\vspace{0.1cm}
		\textbf{OPTIONS:}
		\begin{enumerate}
			\item Removing all punctuation from text
			\item Separating a document into sentences and then into words
			\item Converting all text to lowercase
			\item Removing stop words
		\end{enumerate}
		
		\vspace{0.1cm}
		\textcolor{myblue}{\textbf{ANSWER: B}}
		
		\vspace{0.1cm}
		\textbf{DETAILED TUTORIAL:}
		\begin{itemize}
			\item \textbf{Tokenization process:}
			\begin{enumerate}
				\item \textbf{Sentence tokenization:} Split at periods, question marks, exclamation marks
				\item \textbf{Word tokenization:} Split sentences into words, remove punctuation, convert to lowercase
			\end{enumerate}
		\end{itemize}
		
		\textcolor{myred}{\textbf{EXAM TRAP:}} Tokenization = \textcolor{myblue}{splitting into sentences and words}!
	\end{frame}
	
	\begin{frame}{Q12: Stop Words Removal - Tutorial}
		\fontsize{6.5}{7.5}\selectfont
		\textbf{QUESTION: What are stop words and why are they removed?}
		
		\vspace{0.1cm}
		\textbf{OPTIONS:}
		\begin{enumerate}
			\item Important words that convey meaning
			\item Common words with no informational value (e.g., "the", "a", "also")
			\item All nouns and verbs
			\item All proper nouns
		\end{enumerate}
		
		\vspace{0.1cm}
		\textcolor{myblue}{\textbf{ANSWER: B}}
		
		\vspace{0.1cm}
		\textbf{DETAILED TUTORIAL:}
		\begin{itemize}
			\item \textbf{Stop words definition:}
			\begin{itemize}
				\item Words with little informational value
				\item Examples: "has", "a", "the", "also", "and"
				\item Common in all documents
			\end{itemize}
			\item \textbf{Why remove them:} Add noise, don't help distinguish documents, increase vector size
			\item \textbf{Caveat:} What constitutes a stop word varies by application
		\end{itemize}
		
		\textcolor{myred}{\textbf{EXAM TRAP:}} Stop words = \textcolor{myblue}{common words with no informational value}!
	\end{frame}
	
	\begin{frame}{Q13: Stemming Definition - Tutorial}
		\fontsize{6.5}{7.5}\selectfont
		\textbf{QUESTION: What is stemming in NLP pre-processing?}
		
		\vspace{0.1cm}
		\textbf{OPTIONS:}
		\begin{enumerate}
			\item Removing all words from a document
			\item Replacing words with their core/root forms by cutting prefixes and suffixes
			\item Converting all text to uppercase
			\item Removing all punctuation
		\end{enumerate}
		
		\vspace{0.1cm}
		\textcolor{myblue}{\textbf{ANSWER: B}}
		
		\vspace{0.1cm}
		\textbf{DETAILED TUTORIAL:}
		\begin{itemize}
			\item \textbf{Stemming definition:}
			\begin{itemize}
				\item Reduces words to their root forms
				\item Cuts off prefixes and suffixes
				\item Crude, rule-based approach
			\end{itemize}
			\item \textbf{Examples:} "disappointing" $\to$ "disappoint", "running" $\to$ "run"
			\item \textbf{Limitation:} Can lose nuance, may create non-words
		\end{itemize}
		
		\textcolor{myred}{\textbf{EXAM TRAP:}} Stemming = \textcolor{myblue}{cutting prefixes/suffixes to root form}!
	\end{frame}
	
	\begin{frame}{Q14: Lemmatization Definition - Tutorial}
		\fontsize{6.5}{7.5}\selectfont
		\textbf{QUESTION: What is lemmatization and how does it differ from stemming?}
		
		\vspace{0.1cm}
		\textbf{OPTIONS:}
		\begin{enumerate}
			\item Lemmatization is the same as stemming
			\item Lemmatization finds the dictionary form (lemma) considering word meaning and part of speech
			\item Lemmatization removes stop words
			\item Lemmatization converts text to lowercase
		\end{enumerate}
		
		\vspace{0.1cm}
		\textcolor{myblue}{\textbf{ANSWER: B}}
		
		\vspace{0.1cm}
		\textbf{DETAILED TUTORIAL:}
		\begin{itemize}
			\item \textbf{Lemmatization definition:}
			\begin{itemize}
				\item Finds the dictionary form (lemma) of words
				\item Considers word meaning and part of speech
				\item More sophisticated than stemming
			\end{itemize}
			\item \textbf{Comparison:}
			\begin{itemize}
				\item Stemming: Rule-based, fast, can be inaccurate
				\item Lemmatization: Uses vocabulary, slower, more accurate
			\end{itemize}
		\end{itemize}
		
		\textcolor{myred}{\textbf{EXAM TRAP:}} Lemmatization = \textcolor{myblue}{dictionary form considering meaning and POS}!
	\end{frame}
	
	\begin{frame}{Q15: POS Tagging Purpose - Tutorial}
		\fontsize{6.5}{7.5}\selectfont
		\textbf{QUESTION: What is Part-of-Speech (POS) tagging and why is it useful?}
		
		\vspace{0.1cm}
		\textbf{OPTIONS:}
		\begin{enumerate}
			\item Identifying all verbs in a document
			\item Labelling each word by its class (noun, verb, adjective, proper noun, etc.)
			\item Removing all adjectives from a document
			\item Counting the number of sentences
		\end{enumerate}
		
		\vspace{0.1cm}
		\textcolor{myblue}{\textbf{ANSWER: B}}
		
		\vspace{0.1cm}
		\textbf{DETAILED TUTORIAL:}
		\begin{itemize}
			\item \textbf{POS tagging definition:}
			\begin{itemize}
				\item Labels each word with its grammatical category
				\item Common tags: noun, verb, adjective, proper noun
			\end{itemize}
			\item \textbf{Why it's useful:}
			\begin{itemize}
				\item Handles ambiguity: "Reading" (town) vs "reading" (activity)
				\item Handles case sensitivity: "Polish" (language) vs "polish" (to shine)
			\end{itemize}
		\end{itemize}
		
		\textcolor{myred}{\textbf{EXAM TRAP:}} POS tagging labels \textcolor{myblue}{each word by grammatical class}!
	\end{frame}
	
	\begin{frame}{Q16: Lowercasing Issues - Tutorial}
		\fontsize{6.5}{7.5}\selectfont
		\textbf{QUESTION: What problem can lowercasing cause in NLP?}
		
		\vspace{0.1cm}
		\textbf{OPTIONS:}
		\begin{enumerate}
			\item It makes text harder to read
			\item It can cause inaccuracies by removing case-based distinctions (proper nouns, different meanings)
			\item It doesn't affect NLP
			\item It makes words longer
		\end{enumerate}
		
		\vspace{0.1cm}
		\textcolor{myblue}{\textbf{ANSWER: B}}
		
		\vspace{0.1cm}
		\textbf{DETAILED TUTORIAL:}
		\begin{itemize}
			\item \textbf{Lowercasing problems:}
			\begin{enumerate}
				\item \textbf{Proper nouns:} "Reading" (town) vs "reading" (activity)
				\item \textbf{Meaning loss:} "Polish" (language) vs "polish" (to shine)
			\end{enumerate}
			\item \textbf{Solution:} Apply POS tagging before lowercasing
		\end{itemize}
		
		\textcolor{myred}{\textbf{EXAM TRAP:}} Lowercasing can lose \textcolor{myblue}{proper noun distinctions} - use POS tagging first!
	\end{frame}
	
	\begin{frame}{Q17: Punctuation Removal Effects - Tutorial}
		\fontsize{6.5}{7.5}\selectfont
		\textbf{QUESTION: What is a limitation of removing punctuation?}
		
		\vspace{0.1cm}
		\textbf{OPTIONS:}
		\begin{enumerate}
			\item It doesn't affect NLP
			\item It can cause loss of valuable context (questions, exclamations, sarcasm)
			\item It makes text shorter
			\item It removes all verbs
		\end{enumerate}
		
		\vspace{0.1cm}
		\textcolor{myblue}{\textbf{ANSWER: B}}
		
		\vspace{0.1cm}
		\textbf{DETAILED TUTORIAL:}
		\begin{itemize}
			\item \textbf{Punctuation problems:}
			\begin{enumerate}
				\item \textbf{Questions:} "Is the CEO performing well?" $\to$ may seem positive
				\item \textbf{Exclamations:} May denote humor or sarcasm
			\end{enumerate}
		\end{itemize}
		
		\textcolor{myred}{\textbf{EXAM TRAP:}} Removing punctuation loses \textcolor{myblue}{context and tone} information!
	\end{frame}
	
	\begin{frame}{Q18: Stemming vs Lemmatization Usage - Tutorial}
		\fontsize{6.5}{7.5}\selectfont
		\textbf{QUESTION: When might lemmatization be preferred over stemming?}
		
		\vspace{0.1cm}
		\textbf{OPTIONS:}
		\begin{enumerate}
			\item When speed is the only concern
			\item When accuracy and preserving meaning are important
			\item When working with very short documents
			\item When speed is critical
		\end{enumerate}
		
		\vspace{0.1cm}
		\textcolor{myblue}{\textbf{ANSWER: B}}
		
		\vspace{0.1cm}
		\textbf{DETAILED TUTORIAL:}
		\begin{itemize}
			\item \textbf{When to use lemmatization:}
			\begin{itemize}
				\item Accuracy is critical (financial documents need precise analysis)
				\item Meaning preservation (irregular verbs, context-specific forms)
			\end{itemize}
			\item \textbf{When to use stemming:}
			\begin{itemize}
				\item Speed is critical (very large documents, real-time processing)
			\end{itemize}
			\item \textbf{Tradeoff:} Stemming = faster, less accurate; Lemmatization = slower, more accurate
		\end{itemize}
		
		\textcolor{myred}{\textbf{EXAM TRAP:}} Lemmatization preferred for \textcolor{myblue}{accuracy}, stemming for \textcolor{myblue}{speed}!
	\end{frame}
	
	\begin{frame}{Q19: Vocabulary Building - Tutorial}
		\fontsize{6.5}{7.5}\selectfont
		\textbf{QUESTION: How is the vocabulary for BoW determined?}
		
		\vspace{0.1cm}
		\textbf{OPTIONS:}
		\begin{enumerate}
			\item Only using pre-existing dictionaries
			\item Either imposed exogenously or created from every distinct word in the corpus
			\item Only using the most common words
			\item Only using the rarest words
		\end{enumerate}
		
		\vspace{0.1cm}
		\textcolor{myblue}{\textbf{ANSWER: B}}
		
		\vspace{0.1cm}
		\textbf{DETAILED TUTORIAL:}
		\begin{itemize}
			\item \textbf{Vocabulary creation methods:}
			\begin{enumerate}
				\item \textbf{Exogenous creation:} Created ahead of time, applied to documents
				\item \textbf{Corpus-based creation:} Every distinct word in all documents
			\end{enumerate}
		\end{itemize}
		
		\textcolor{myred}{\textbf{EXAM TRAP:}} Vocabulary = \textcolor{myblue}{pre-defined OR from corpus words}!
	\end{frame}
	
	\begin{frame}{Q20: Bag of Words Definition - Tutorial}
		\fontsize{6.5}{7.5}\selectfont
		\textbf{QUESTION: What is the Bag of Words (BoW) approach?}
		
		\vspace{0.1cm}
		\textbf{OPTIONS:}
		\begin{enumerate}
			\item Analyzing words in their original order
			\item Treating each word in a document independently, ignoring order
			\item Counting only the longest words
			\item Removing all words from documents
		\end{enumerate}
		
		\vspace{0.1cm}
		\textcolor{myblue}{\textbf{ANSWER: B}}
		
		\vspace{0.1cm}
		\textbf{DETAILED TUTORIAL:}
		\begin{itemize}
			\item \textbf{BoW definition:}
			\begin{itemize}
				\item Treats each word distinctly and equally
				\item Position and order are ignored
				\item Counts occurrences in each document
			\end{itemize}
			\item \textbf{Limitations:}
			\begin{itemize}
				\item Loses context
				\item Ignores negation words
				\item Misses phrase meanings
			\end{itemize}
		\end{itemize}
		
		\textcolor{myred}{\textbf{EXAM TRAP:}} BoW ignores \textcolor{myblue}{word order and position}!
	\end{frame}
	
	\begin{frame}{Q21: Binary vs Count BoW - Tutorial}
		\fontsize{6.5}{7.5}\selectfont
		\textbf{QUESTION: What is the difference between binary and count Bag of Words?}
		
		\vspace{0.1cm}
		\textbf{OPTIONS:}
		\begin{enumerate}
			\item Binary uses 1/0, count uses frequency counts
			\item Binary is more accurate than count
			\item Count ignores word frequencies
			\item There is no difference
		\end{enumerate}
		
		\vspace{0.1cm}
		\textcolor{myblue}{\textbf{ANSWER: A}}
		
		\vspace{0.1cm}
		\textbf{DETAILED TUTORIAL:}
		\begin{itemize}
			\item \textbf{Binary BoW:}
			\begin{itemize}
				\item 1 if word appears, 0 if not
				\item Presence/absence only
			\end{itemize}
			\item \textbf{Count BoW:}
			\begin{itemize}
				\item Counts how many times each word appears
				\item Non-negative integers
			\end{itemize}
		\end{itemize}
		
		\textcolor{myred}{\textbf{EXAM TRAP:}} Binary = \textcolor{myblue}{presence/absence}, Count = \textcolor{myblue}{frequency counts}!
	\end{frame}
	
	\begin{frame}{Q22: Document Term Matrix - Tutorial}
		\fontsize{6.5}{7.5}\selectfont
		\textbf{QUESTION: What is a Document Term Matrix (DTM)?}
		
		\vspace{0.1cm}
		\textbf{OPTIONS:}
		\begin{enumerate}
			\item A matrix of document titles
			\item A collection of document vectors with rows as documents and columns as terms
			\item A matrix of document dates
			\item A matrix of document authors
		\end{enumerate}
		
		\vspace{0.1cm}
		\textcolor{myblue}{\textbf{ANSWER: B}}
		
		\vspace{0.1cm}
		\textbf{DETAILED TUTORIAL:}
		\begin{itemize}
			\item \textbf{DTM definition:}
			\begin{itemize}
				\item Rows = documents
				\item Columns = terms/words
				\item Entries = word counts or frequencies
			\end{itemize}
			\item \textbf{Characteristics:}
			\begin{itemize}
				\item Sparse (many zeros)
				\item Large (many columns)
			\end{itemize}
		\end{itemize}
		
		\textcolor{myred}{\textbf{EXAM TRAP:}} DTM = \textcolor{myblue}{documents as rows, terms as columns}!
	\end{frame}
	
	\begin{frame}{Q23: Sparse Vectors Problem - Tutorial}
		\fontsize{6.5}{7.5}\selectfont
		\textbf{QUESTION: Why are document word vectors described as sparse?}
		
		\vspace{0.1cm}
		\textbf{OPTIONS:}
		\begin{enumerate}
			\item They contain many non-zero values
			\item They contain many zeros
			\item They are all the same length
			\item They are very small
		\end{enumerate}
		
		\vspace{0.1cm}
		\textcolor{myblue}{\textbf{ANSWER: B}}
		
		\vspace{0.1cm}
		\textbf{DETAILED TUTORIAL:}
		\begin{itemize}
			\item \textbf{Sparsity definition:}
			\begin{itemize}
				\item Most words don't appear in each document
				\item Vector entries are mostly zero
			\end{itemize}
			\item \textbf{Problems:}
			\begin{itemize}
				\item Computational inefficiency
				\item Curse of dimensionality
			\end{itemize}
		\end{itemize}
		
		\textcolor{myred}{\textbf{EXAM TRAP:}} Sparse = \textcolor{myblue}{many zeros} - makes modeling inefficient!
	\end{frame}
	
	\begin{frame}{Q24: Vector Normalization - Tutorial}
		\fontsize{6.5}{7.5}\selectfont
		\textbf{QUESTION: Why is vector normalization needed in NLP?}
		
		\vspace{0.1cm}
		\textbf{OPTIONS:}
		\begin{enumerate}
			\item To make vectors longer
			\item To prevent repeated words from skewing analysis
			\item To remove all words
			\item To make vectors shorter
		\end{enumerate}
		
		\vspace{0.1cm}
		\textcolor{myblue}{\textbf{ANSWER: B}}
		
		\vspace{0.1cm}
		\textbf{DETAILED TUTORIAL:}
		\begin{itemize}
			\item \textbf{Why normalize:}
			\begin{itemize}
				\item Repeated words can skew analysis
				\item Similar to outliers in econometrics
			\end{itemize}
			\item \textbf{L2 normalization:}
			\begin{itemize}
				\item Divide by the L2 norm
				\item Scales elements between 0 and 1
				\item All documents become unit length
			\end{itemize}
		\end{itemize}
		
		\textcolor{myred}{\textbf{EXAM TRAP:}} L2 normalization = \textcolor{myblue}{scale to unit length}!
	\end{frame}
	
	\begin{frame}{Q25: L2 Norm Formula - Tutorial}
		\fontsize{6.5}{7.5}\selectfont
		\textbf{QUESTION: What is the formula for the L2 norm of a vector?}
		
		\vspace{0.1cm}
		\textbf{OPTIONS:}
		\begin{enumerate}
			\item $||X||_2 = \sum x_i$
			\item $||X||_2 = \sqrt{\sum x_i^2}$
			\item $||X||_2 = \sum x_i^2$
			\item $||X||_2 = \max(x_i)$
		\end{enumerate}
		
		\vspace{0.1cm}
		\textcolor{myblue}{\textbf{ANSWER: B}}
		
		\vspace{0.1cm}
		\textbf{DETAILED TUTORIAL:}
		\begin{itemize}
			\item \textbf{L2 norm formula:}
			\begin{itemize}
				\item $||X||_2 = \sqrt{\sum_{i=1}^n x_i^2}$
				\item Also called Euclidean norm
			\end{itemize}
			\item \textbf{Application in NLP:}
			\begin{itemize}
				\item Normalize document vectors
				\item Each vector becomes unit length
			\end{itemize}
		\end{itemize}
		
		\textcolor{myred}{\textbf{EXAM TRAP:}} L2 norm = \textcolor{myblue}{square root of sum of squared elements}!
	\end{frame}
	
	% ============================================================
	% SECTION 3: DICTIONARY APPROACH - QUESTIONS 26-40
	% ============================================================
	
	\begin{frame}{Q26: Dictionary Approach Definition - Tutorial}
		\fontsize{6.5}{7.5}\selectfont
		\textbf{QUESTION: What is the dictionary approach in NLP?}
		
		\vspace{0.1cm}
		\textbf{OPTIONS:}
		\begin{enumerate}
			\item A machine learning algorithm
			\item A pre-defined list of words sharing a common characteristic used for text analysis
			\item A method for removing stop words
			\item A way to tokenize documents
		\end{enumerate}
		
		\vspace{0.1cm}
		\textcolor{myblue}{\textbf{ANSWER: B}}
		
		\vspace{0.1cm}
		\textbf{DETAILED TUTORIAL:}
		\begin{itemize}
			\item \textbf{Dictionary approach definition:}
			\begin{itemize}
				\item Pre-defined list of words
				\item Words share a common characteristic
				\item Used to assess sentiment or topic
			\end{itemize}
			\item \textbf{Common application:}
			\begin{itemize}
				\item Sentiment analysis
				\item Net sentiment = proportion positive - proportion negative
			\end{itemize}
		\end{itemize}
		
		\textcolor{myred}{\textbf{EXAM TRAP:}} Dictionary = \textcolor{myblue}{pre-defined word lists with common characteristics}!
	\end{frame}
	
	\begin{frame}{Q27: Dictionary Approach Example - Tutorial}
		\fontsize{6.5}{7.5}\selectfont
		\textbf{QUESTION: In the newswire example, how is sentiment determined?}
		
		\vspace{0.1cm}
		\textbf{OPTIONS:}
		\begin{enumerate}
			\item Count only positive words
			\item Count only negative words
			\item Count positive and negative words, subtract negatives from positives
			\item Count all words equally
		\end{enumerate}
		
		\vspace{0.1cm}
		\textcolor{myblue}{\textbf{ANSWER: C}}
		
		\vspace{0.1cm}
		\textbf{DETAILED TUTORIAL:}
		\begin{itemize}
			\item \textbf{Sentiment calculation:}
			\begin{itemize}
				\item Count positive words: rise, grow, resolve, relief
				\item Count negative words: disappoint, worry, concern, decline
				\item Net sentiment = positive count - negative count
			\end{itemize}
			\item \textbf{Challenges:}
			\begin{itemize}
				\item Counterfactual sentences
				\item Complex sentence structure
				\item Negation words can reverse meaning
			\end{itemize}
		\end{itemize}
		
		\textcolor{myred}{\textbf{EXAM TRAP:}} Net sentiment = \textcolor{myblue}{positive count - negative count}!
	\end{frame}
	
	\begin{frame}{Q28: Loughran-McDonald Dictionary - Tutorial}
		\fontsize{6.5}{7.5}\selectfont
		\textbf{QUESTION: What is the Loughran-McDonald dictionary known for?}
		
		\vspace{0.1cm}
		\textbf{OPTIONS:}
		\begin{enumerate}
			\item Being the largest dictionary
			\item Being specifically designed for financial text analysis
			\item Being the oldest dictionary
			\item Only containing positive words
		\end{enumerate}
		
		\vspace{0.1cm}
		\textcolor{myblue}{\textbf{ANSWER: B}}
		
		\vspace{0.1cm}
		\textbf{DETAILED TUTORIAL:}
		\begin{itemize}
			\item \textbf{Loughran-McDonald dictionary:}
			\begin{itemize}
				\item Developed for financial text analysis
				\item Based on extensive examination of 10-K filings
				\item Categories: positive, negative, uncertain, litigious, strong modal, weak modal
			\end{itemize}
			\item \textbf{Accuracy:} Around 75\% on Thomson Reuters news articles
		\end{itemize}
		
		\textcolor{myred}{\textbf{EXAM TRAP:}} Loughran-McDonald = \textcolor{myblue}{financial text dictionary} - 75\% accuracy!
	\end{frame}
	
	\begin{frame}{Q29: Dictionary Approach Advantages - Tutorial}
		\fontsize{6.5}{7.5}\selectfont
		\textbf{QUESTION: What are the primary advantages of the dictionary approach?}
		
		\vspace{0.1cm}
		\textbf{OPTIONS:}
		\begin{enumerate}
			\item It's always 100\% accurate
			\item Transparency, ease and speed of implementation
			\item It works for all types of text
			\item It requires no pre-processing
		\end{enumerate}
		
		\vspace{0.1cm}
		\textcolor{myblue}{\textbf{ANSWER: B}}
		
		\vspace{0.1cm}
		\textbf{DETAILED TUTORIAL:}
		\begin{itemize}
			\item \textbf{Key advantages:}
			\begin{enumerate}
				\item \textbf{Transparency:} Word lists are explicit and clear
				\item \textbf{Ease of implementation:} Simple word counting
				\item \textbf{Comparability:} Same word lists across corpora
			\end{enumerate}
		\end{itemize}
		
		\textcolor{myred}{\textbf{EXAM TRAP:}} Advantages = \textcolor{myblue}{transparency, ease, speed, comparability}!
	\end{frame}
	
	\begin{frame}{Q30: Dictionary Approach Limitations - Tutorial}
		\fontsize{6.5}{7.5}\selectfont
		\textbf{QUESTION: What is a key limitation of the dictionary approach?}
		
		\vspace{0.1cm}
		\textbf{OPTIONS:}
		\begin{enumerate}
			\item It's too fast
			\item Dictionaries are only as good as the word lists they are built on
			\item It always overestimates sentiment
			\item It always underestimates sentiment
		\end{enumerate}
		
		\vspace{0.1cm}
		\textcolor{myblue}{\textbf{ANSWER: B}}
		
		\vspace{0.1cm}
		\textbf{DETAILED TUTORIAL:}
		\begin{itemize}
			\item \textbf{Key limitations:}
			\begin{enumerate}
				\item \textbf{Quality dependence:} Only as good as word lists
				\item \textbf{Context limitation:} Designed for formal documents
				\item \textbf{Narrow scope:} Limited subjects available
				\item \textbf{Negation issues:} "not good" misclassified
			\end{enumerate}
		\end{itemize}
		
		\textcolor{myred}{\textbf{EXAM TRAP:}} Dictionary = \textcolor{myblue}{only as good as word lists} - has many limitations!
	\end{frame}
	
	\begin{frame}{Q31: n-gram Definition - Tutorial}
		\fontsize{6.5}{7.5}\selectfont
		\textbf{QUESTION: What is an n-gram in NLP?}
		
		\vspace{0.1cm}
		\textbf{OPTIONS:}
		\begin{enumerate}
			\item A single word
			\item A sequence of n contiguous words treated as a single unit
			\item A type of machine learning algorithm
			\item A pre-processing technique
		\end{enumerate}
		
		\vspace{0.1cm}
		\textcolor{myblue}{\textbf{ANSWER: B}}
		
		\vspace{0.1cm}
		\textbf{DETAILED TUTORIAL:}
		\begin{itemize}
			\item \textbf{n-gram definition:}
			\begin{itemize}
				\item n = 1: uni-gram (single word)
				\item n = 2: bi-gram (two words)
				\item n = 3: tri-gram (three words)
			\end{itemize}
			\item \textbf{Example:} "credit spreads narrowed today" $\to$ bi-grams: ["credit spreads", "spreads narrowed", "narrowed today"]
			\item \textbf{Purpose:} Capture phrase meanings, handle negation
		\end{itemize}
		
		\textcolor{myred}{\textbf{EXAM TRAP:}} n-gram = \textcolor{myblue}{n contiguous words as one unit}!
	\end{frame}
	
	\begin{frame}{Q32: Why Use n-grams? - Tutorial}
		\fontsize{6.5}{7.5}\selectfont
		\textbf{QUESTION: Why are n-grams needed beyond basic BoW?}
		
		\vspace{0.1cm}
		\textbf{OPTIONS:}
		\begin{enumerate}
			\item To make analysis faster
			\item Because some words have specific meaning when placed together
			\item Because basic BoW is too accurate
			\item To count words more precisely
		\end{enumerate}
		
		\vspace{0.1cm}
		\textcolor{myblue}{\textbf{ANSWER: B}}
		
		\vspace{0.1cm}
		\textbf{DETAILED TUTORIAL:}
		\begin{itemize}
			\item \textbf{Limitations of basic BoW:}
			\begin{enumerate}
				\item \textbf{Loses phrase meaning:} "red herring" $\neq$ red + herring
				\item \textbf{Negation words:} "not good" $\neq$ not + good
			\end{enumerate}
			\item \textbf{How n-grams help:}
			\begin{itemize}
				\item Preserve phrase meaning
				\item Handle negations correctly
				\item Capture context
			\end{itemize}
		\end{itemize}
		
		\textcolor{myred}{\textbf{EXAM TRAP:}} n-grams capture \textcolor{myblue}{phrase meaning and negation context}!
	\end{frame}
	
	\begin{frame}{Q33: Negation Handling with n-grams - Tutorial}
		\fontsize{6.5}{7.5}\selectfont
		\textbf{QUESTION: How can n-grams help with negation words?}
		
		\vspace{0.1cm}
		\textbf{OPTIONS:}
		\begin{enumerate}
			\item By removing all negation words
			\item By treating "not" + following word as a single unit
			\item By ignoring negations
			\item By making negations positive
		\end{enumerate}
		
		\vspace{0.1cm}
		\textcolor{myblue}{\textbf{ANSWER: B}}
		
		\vspace{0.1cm}
		\textbf{DETAILED TUTORIAL:}
		\begin{itemize}
			\item \textbf{Negation problem:}
			\begin{itemize}
				\item "not good" $\to$ should be negative
				\item Basic BoW sees "not" and "good" separately
			\end{itemize}
			\item \textbf{n-gram solution:}
			\begin{itemize}
				\item Bi-grams capture "not good" as one unit
				\item Preserves negation meaning
			\end{itemize}
		\end{itemize}
		
		\textcolor{myred}{\textbf{EXAM TRAP:}} n-grams pair \textcolor{myblue}{negation + following word} as one unit!
	\end{frame}
	
	\begin{frame}{Q34: TF-IDF Definition - Tutorial}
		\fontsize{6.5}{7.5}\selectfont
		\textbf{QUESTION: What does TF-IDF stand for and what does it measure?}
		
		\vspace{0.1cm}
		\textbf{OPTIONS:}
		\begin{enumerate}
			\item Term Frequency - Inverse Document Frequency; weights words by importance
			\item Total Frequency - Integrated Document Frequency
			\item Text Frequency - Inverse Data Frequency
			\item Term Frequency - Inverse Data Frequency
		\end{enumerate}
		
		\vspace{0.1cm}
		\textcolor{myblue}{\textbf{ANSWER: A}}
		
		\vspace{0.1cm}
		\textbf{DETAILED TUTORIAL:}
		\begin{itemize}
			\item \textbf{TF-IDF definition:}
			\begin{itemize}
				\item \textbf{TF (Term Frequency):} How often word appears in document
				\item \textbf{IDF (Inverse Document Frequency):} How rare the word is across documents
				\item \textbf{TF-IDF:} TF $\times$ IDF = word importance score
			\end{itemize}
		\end{itemize}
		
		\textcolor{myred}{\textbf{EXAM TRAP:}} TF-IDF = \textcolor{myblue}{Term Frequency $\times$ Inverse Document Frequency}!
	\end{frame}
	
	\begin{frame}{Q35: Term Frequency Formula - Tutorial}
		\fontsize{6.5}{7.5}\selectfont
		\textbf{QUESTION: What is the formula for Term Frequency (TF)?}
		
		\vspace{0.1cm}
		\textbf{OPTIONS:}
		\begin{enumerate}
			\item $TF_{i,j} = w_{i,j} \times |v_j|$
			\item $TF_{i,j} = w_{i,j} / |v_j|$
			\item $TF_{i,j} = |v_j| / w_{i,j}$
			\item $TF_{i,j} = \log(w_{i,j})$
		\end{enumerate}
		
		\vspace{0.1cm}
		\textcolor{myblue}{\textbf{ANSWER: B}}
		
		\vspace{0.1cm}
		\textbf{DETAILED TUTORIAL:}
		\begin{itemize}
			\item \textbf{TF Formula:}
			\begin{itemize}
				\item $TF_{i,j} = \frac{w_{i,j}}{|v_j|}$
				\item $w_{i,j}$: number of times term i appears in document j
				\item $|v_j|$: total number of terms in document j
			\end{itemize}
		\end{itemize}
		
		\textcolor{myred}{\textbf{EXAM TRAP:}} TF = \textcolor{myblue}{word count / total words in document}!
	\end{frame}
	
	\begin{frame}{Q36: Inverse Document Frequency Formula - Tutorial}
		\fontsize{6.5}{7.5}\selectfont
		\textbf{QUESTION: What is the formula for Inverse Document Frequency (IDF)?}
		
		\vspace{0.1cm}
		\textbf{OPTIONS:}
		\begin{enumerate}
			\item $IDF_i = \ln(D / \sum d_{i,j})$
			\item $IDF_i = D / \sum d_{i,j}$
			\item $IDF_i = \sum d_{i,j} / D$
			\item $IDF_i = \ln(\sum d_{i,j} / D)$
		\end{enumerate}
		
		\vspace{0.1cm}
		\textcolor{myblue}{\textbf{ANSWER: A}}
		
		\vspace{0.1cm}
		\textbf{DETAILED TUTORIAL:}
		\begin{itemize}
			\item \textbf{IDF Formula:}
			\begin{itemize}
				\item $IDF_i = \ln(D / \sum_{j} d_{i,j})$
				\item D = total number of documents
				\item $\sum d_{i,j}$ = number of documents containing term i
			\end{itemize}
			\item \textbf{Interpretation:}
			\begin{itemize}
				\item Rare words: High IDF
				\item Common words: Low IDF
				\item Words in all documents: IDF = 0
			\end{itemize}
		\end{itemize}
		
		\textcolor{myred}{\textbf{EXAM TRAP:}} IDF = \textcolor{myblue}{ln(total documents / documents containing word)}!
	\end{frame}
	
	\begin{frame}{Q37: Zero TF-IDF Value - Tutorial}
		\fontsize{6.5}{7.5}\selectfont
		\textbf{QUESTION: Why do some words have zero TF-IDF values?}
		
		\vspace{0.1cm}
		\textbf{OPTIONS:}
		\begin{enumerate}
			\item Because they are not in the vocabulary
			\item Because they appear in every document (no discriminatory power)
			\item Because they are too long
			\item Because they are too short
		\end{enumerate}
		
		\vspace{0.1cm}
		\textcolor{myblue}{\textbf{ANSWER: B}}
		
		\vspace{0.1cm}
		\textbf{DETAILED TUTORIAL:}
		\begin{itemize}
			\item \textbf{Zero IDF explanation:}
			\begin{itemize}
				\item IDF = ln(D / $\sum d_{i,j}$)
				\item If word appears in all documents: $\sum d_{i,j} = D$
				\item IDF = ln(D/D) = ln(1) = 0
			\end{itemize}
		\end{itemize}
		
		\textcolor{myred}{\textbf{EXAM TRAP:}} Zero TF-IDF = \textcolor{myblue}{word appears in all documents} - no discriminatory power!
	\end{frame}
	
	\begin{frame}{Q38: TF-IDF Variation by Document - Tutorial}
		\fontsize{6.5}{7.5}\selectfont
		\textbf{QUESTION: Why does the same word have different TF-IDF scores across documents?}
		
		\vspace{0.1cm}
		\textbf{OPTIONS:}
		\begin{enumerate}
			\item IDF varies by document
			\item TF varies by document (due to different lengths and frequencies)
			\item TF-IDF is always constant
			\item Vocabulary changes
		\end{enumerate}
		
		\vspace{0.1cm}
		\textcolor{myblue}{\textbf{ANSWER: B}}
		
		\vspace{0.1cm}
		\textbf{DETAILED TUTORIAL:}
		\begin{itemize}
			\item \textbf{Why TF varies:}
			\begin{enumerate}
				\item \textbf{Document length:} Longer documents have smaller TFs
				\item \textbf{Frequency:} Word may appear multiple times
			\end{enumerate}
			\item \textbf{IDF is constant:} Depends only on corpus
		\end{itemize}
		
		\textcolor{myred}{\textbf{EXAM TRAP:}} IDF is \textcolor{myblue}{constant across documents}; TF varies $\Rightarrow$ TF-IDF varies!
	\end{frame}
	
	\begin{frame}{Q39: TF-IDF Variations - Tutorial}
		\fontsize{6.5}{7.5}\selectfont
		\textbf{QUESTION: What are common variations in TF-IDF implementation?}
		
		\vspace{0.1cm}
		\textbf{OPTIONS:}
		\begin{enumerate}
			\item Only one formula exists
			\item Variations include binary TF, log normalization, and different IDF formulas
			\item TF-IDF is always the same
			\item TF-IDF cannot be modified
		\end{enumerate}
		
		\vspace{0.1cm}
		\textcolor{myblue}{\textbf{ANSWER: B}}
		
		\vspace{0.1cm}
		\textbf{DETAILED TUTORIAL:}
		\begin{itemize}
			\item \textbf{Common variations:}
			\begin{enumerate}
				\item \textbf{TF variations:} Binary TF, log normalization, raw count
				\item \textbf{IDF variations:} Smooth IDF, probabilistic IDF
			\end{enumerate}
		\end{itemize}
		
		\textcolor{myred}{\textbf{EXAM TRAP:}} TF-IDF has \textcolor{myblue}{many implementation variations} - be aware!
	\end{frame}
	
	\begin{frame}{Q40: TF-IDF Example Interpretation - Tutorial}
		\fontsize{6.5}{7.5}\selectfont
		\textbf{QUESTION: In the TF-IDF example, which word has the highest score and why?}
		
		\vspace{0.1cm}
		\textbf{OPTIONS:}
		\begin{enumerate}
			\item "this" - because it appears in all documents
			\item "fed" - because it appears in only one short document
			\item "down" - because it appears frequently
			\item "stock" - because it appears in two documents
		\end{enumerate}
		
		\vspace{0.1cm}
		\textcolor{myblue}{\textbf{ANSWER: B}}
		
		\vspace{0.1cm}
		\textbf{DETAILED TUTORIAL:}
		\begin{itemize}
			\item \textbf{Highest scores:} "fed" = 0.28, "hikes" = 0.28, "rates" = 0.28
			\item \textbf{Why they're highest:}
			\begin{itemize}
				\item Appear in only 1 out of 4 documents
				\item IDF = ln(4/1) = 1.39 (high)
				\item Document length = 5 (short)
				\item TF = 1/5 = 0.20
				\item TF-IDF = 0.20 $\times$ 1.39 = 0.28
			\end{itemize}
		\end{itemize}
		
		\textcolor{myred}{\textbf{EXAM TRAP:}} Rare words in short documents = \textcolor{myblue}{highest TF-IDF}!
	\end{frame}
	
	% ============================================================
	% SECTION 4: NAÏVE BAYES CLASSIFIER - QUESTIONS 41-60
	% ============================================================
	
	\begin{frame}{Q41: Naïve Bayes Definition - Tutorial}
		\fontsize{6.5}{7.5}\selectfont
		\textbf{QUESTION: What is the Naïve Bayes classifier?}
		
		\vspace{0.1cm}
		\textbf{OPTIONS:}
		\begin{enumerate}
			\item A regression algorithm
			\item A classification algorithm based on Bayes' theorem with independence assumption
			\item A clustering algorithm
			\item A dimensionality reduction technique
		\end{enumerate}
		
		\vspace{0.1cm}
		\textcolor{myblue}{\textbf{ANSWER: B}}
		
		\vspace{0.1cm}
		\textbf{DETAILED TUTORIAL:}
		\begin{itemize}
			\item \textbf{Naïve Bayes definition:}
			\begin{itemize}
				\item Based on Bayes' theorem
				\item Assumes features are independent (naïve assumption)
				\item Used for text classification
			\end{itemize}
			\item \textbf{Why "naïve":} Assumes word independence (unrealistic but works)
		\end{itemize}
		
		\textcolor{myred}{\textbf{EXAM TRAP:}} Naïve = \textcolor{myblue}{independence assumption} - words treated independently!
	\end{frame}
	
	\begin{frame}{Q42: Bayes' Theorem - Tutorial}
		\fontsize{6.5}{7.5}\selectfont
		\textbf{QUESTION: What is Bayes' theorem in the context of classification?}
		
		\vspace{0.1cm}
		\textbf{OPTIONS:}
		\begin{enumerate}
			\item $P(c|d) = P(d|c)P(c) / P(d)$
			\item $P(c|d) = P(d|c) + P(c)$
			\item $P(c|d) = P(c) / P(d)$
			\item $P(c|d) = P(d) / P(c)$
		\end{enumerate}
		
		\vspace{0.1cm}
		\textcolor{myblue}{\textbf{ANSWER: A}}
		
		\vspace{0.1cm}
		\textbf{DETAILED TUTORIAL:}
		\begin{itemize}
			\item \textbf{Bayes' Theorem:}
			\begin{itemize}
				\item $P(c|d) = \frac{P(d|c) \times P(c)}{P(d)}$
				\item $P(c|d)$: Posterior (probability of class given document)
				\item $P(d|c)$: Likelihood (probability of document given class)
				\item $P(c)$: Prior (unconditional probability of class)
			\end{itemize}
		\end{itemize}
		
		\textcolor{myred}{\textbf{EXAM TRAP:}} Bayes = \textcolor{myblue}{posterior = likelihood $\times$ prior / evidence}!
	\end{frame}
	
	\begin{frame}{Q43: Naïve Independence Assumption - Tutorial}
		\fontsize{6.5}{7.5}\selectfont
		\textbf{QUESTION: What is the independence assumption in Naïve Bayes?}
		
		\vspace{0.1cm}
		\textbf{OPTIONS:}
		\begin{enumerate}
			\item Documents are independent of each other
			\item Each word in a document is independent of all other words
			\item Classes are independent of documents
			\item All documents have the same length
		\end{enumerate}
		
		\vspace{0.1cm}
		\textcolor{myblue}{\textbf{ANSWER: B}}
		
		\vspace{0.1cm}
		\textbf{DETAILED TUTORIAL:}
		\begin{itemize}
			\item \textbf{Independence assumption:}
			\begin{itemize}
				\item $P(w_1, w_2, ..., w_n | c) = P(w_1|c) \times P(w_2|c) \times ...$
				\item Words are conditionally independent given class
			\end{itemize}
		\end{itemize}
		
		\textcolor{myred}{\textbf{EXAM TRAP:}} Independence = \textcolor{myblue}{words treated as independent} - unrealistic but works!
	\end{frame}
	
	\begin{frame}{Q44: Prior Probability - Tutorial}
		\fontsize{6.5}{7.5}\selectfont
		\textbf{QUESTION: How is prior probability $P(c)$ calculated in Naïve Bayes?}
		
		\vspace{0.1cm}
		\textbf{OPTIONS:}
		\begin{enumerate}
			\item $P(c) = n(c) / N$ (documents in class / total documents)
			\item $P(c) = N / n(c)$
			\item $P(c) = n(c) \times N$
			\item $P(c) = 1 / N$
		\end{enumerate}
		
		\vspace{0.1cm}
		\textcolor{myblue}{\textbf{ANSWER: A}}
		
		\vspace{0.1cm}
		\textbf{DETAILED TUTORIAL:}
		\begin{itemize}
			\item \textbf{Prior probability formula:}
			\begin{itemize}
				\item $P(c_i) = \frac{n(c_i)}{N}$
				\item $n(c_i)$: number of documents labeled as class i
				\item $N$: total number of documents
			\end{itemize}
		\end{itemize}
		
		\textcolor{myred}{\textbf{EXAM TRAP:}} Prior = \textcolor{myblue}{class count / total documents}!
	\end{frame}
	
	\begin{frame}{Q45: Likelihood Calculation - Tutorial}
		\fontsize{6.5}{7.5}\selectfont
		\textbf{QUESTION: How is the likelihood $P(w|c)$ calculated in Naïve Bayes?}
		
		\vspace{0.1cm}
		\textbf{OPTIONS:}
		\begin{enumerate}
			\item Count of word occurrences in class / total words in class
			\item Count of documents in class / total documents
			\item Word count / total words in corpus
			\item Always 0.5
		\end{enumerate}
		
		\vspace{0.1cm}
		\textcolor{myblue}{\textbf{ANSWER: A}}
		
		\vspace{0.1cm}
		\textbf{DETAILED TUTORIAL:}
		\begin{itemize}
			\item \textbf{Likelihood calculation:}
			\begin{itemize}
				\item $P(w_i | c_j) = \frac{\text{count of } w_i \text{ in class } c_j}{\text{total words in class } c_j}$
				\item Estimated from labeled training data
			\end{itemize}
		\end{itemize}
		
		\textcolor{myred}{\textbf{EXAM TRAP:}} Likelihood = \textcolor{myblue}{word count in class / total words in class}!
	\end{frame}
	
	\begin{frame}{Q46: Posterior Probability - Tutorial}
		\fontsize{6.5}{7.5}\selectfont
		\textbf{QUESTION: What is the posterior probability in Naïve Bayes?}
		
		\vspace{0.1cm}
		\textbf{OPTIONS:}
		\begin{enumerate}
			\item Probability of a document given its class
			\item Probability of a class given the document
			\item Probability of a word given its class
			\item The unconditional class probability
		\end{enumerate}
		
		\vspace{0.1cm}
		\textcolor{myblue}{\textbf{ANSWER: B}}
		
		\vspace{0.1cm}
		\textbf{DETAILED TUTORIAL:}
		\begin{itemize}
			\item \textbf{Posterior definition:}
			\begin{itemize}
				\item $P(c|d)$
				\item Probability that document d belongs to class c
				\item Updated after seeing the document
			\end{itemize}
		\end{itemize}
		
		\textcolor{myred}{\textbf{EXAM TRAP:}} Posterior = \textcolor{myblue}{probability of class given document}!
	\end{frame}
	
	\begin{frame}{Q47: Maximum A Posteriori - Tutorial}
		\fontsize{6.5}{7.5}\selectfont
		\textbf{QUESTION: What is Maximum A Posteriori (MAP) classification?}
		
		\vspace{0.1cm}
		\textbf{OPTIONS:}
		\begin{enumerate}
			\item Assign document to the class with the highest posterior probability
			\item Assign document to the class with the highest prior probability
			\item Assign document to the class with the highest likelihood
			\item Random classification
		\end{enumerate}
		
		\vspace{0.1cm}
		\textcolor{myblue}{\textbf{ANSWER: A}}
		
		\vspace{0.1cm}
		\textbf{DETAILED TUTORIAL:}
		\begin{itemize}
			\item \textbf{MAP definition:}
			\begin{itemize}
				\item $\hat{c} = \arg\max_{c \in C} P(c|d)$
				\item Choose class that maximizes posterior
			\end{itemize}
		\end{itemize}
		
		\textcolor{myred}{\textbf{EXAM TRAP:}} MAP = \textcolor{myblue}{class with highest posterior probability}!
	\end{frame}
	
	\begin{frame}{Q48: Zero Probability Problem - Tutorial}
		\fontsize{6.5}{7.5}\selectfont
		\textbf{QUESTION: What is the zero probability problem in Naïve Bayes?}
		
		\vspace{0.1cm}
		\textbf{OPTIONS:}
		\begin{enumerate}
			\item When a word appears in all documents
			\item When a word does not appear in a class in the training data, making the class probability zero
			\item When all words have zero frequency
			\item When no documents exist
		\end{enumerate}
		
		\vspace{0.1cm}
		\textcolor{myblue}{\textbf{ANSWER: B}}
		
		\vspace{0.1cm}
		\textbf{DETAILED TUTORIAL:}
		\begin{itemize}
			\item \textbf{Zero probability problem:}
			\begin{itemize}
				\item Word doesn't appear in training data for a class
				\item $P(w|c) = 0$
				\item Product becomes zero: $\prod P(w_i|c) = 0$
			\end{itemize}
		\end{itemize}
		
		\textcolor{myred}{\textbf{EXAM TRAP:}} Zero probability = \textcolor{myblue}{word missing in training class} - need smoothing!
	\end{frame}
	
	\begin{frame}{Q49: Laplace Smoothing - Tutorial}
		\fontsize{6.5}{7.5}\selectfont
		\textbf{QUESTION: What is Laplace smoothing in Naïve Bayes?}
		
		\vspace{0.1cm}
		\textbf{OPTIONS:}
		\begin{enumerate}
			\item A method to remove stop words
			\item Adding 1 to all counts to avoid zero probabilities
			\item A technique to stem words
			\item A method to normalize vectors
		\end{enumerate}
		
		\vspace{0.1cm}
		\textcolor{myblue}{\textbf{ANSWER: B}}
		
		\vspace{0.1cm}
		\textbf{DETAILED TUTORIAL:}
		\begin{itemize}
			\item \textbf{Laplace smoothing:}
			\begin{itemize}
				\item Adds 1 to all counts
				\item Ensures all probabilities are non-zero
				\item $P(w|c) = \frac{\text{count}(w,c) + 1}{\text{total words in c} + |V|}$
			\end{itemize}
		\end{itemize}
		
		\textcolor{myred}{\textbf{EXAM TRAP:}} Laplace smoothing = \textcolor{myblue}{add 1 to all counts to avoid zeros}!
	\end{frame}
	
	\begin{frame}{Q50: Smoothing Impact on Probabilities - Tutorial}
		\fontsize{6.5}{7.5}\selectfont
		\textbf{QUESTION: How does smoothing affect the estimated probabilities?}
		
		\vspace{0.1cm}
		\textbf{OPTIONS:}
		\begin{enumerate}
			\item It makes probabilities more extreme
			\item It pushes probabilities toward a uniform distribution
			\item It has no effect on probabilities
			\item It makes probabilities negative
		\end{enumerate}
		
		\vspace{0.1cm}
		\textcolor{myblue}{\textbf{ANSWER: B}}
		
		\vspace{0.1cm}
		\textbf{DETAILED TUTORIAL:}
		\begin{itemize}
			\item \textbf{Effect of smoothing:}
			\begin{itemize}
				\item Prevents zero probabilities
				\item Shrinks probabilities toward uniform
				\item Similar to regularization (Lasso/Ridge)
			\end{itemize}
		\end{itemize}
		
		\textcolor{myred}{\textbf{EXAM TRAP:}} Smoothing = \textcolor{myblue}{shrinks toward uniform distribution} - similar to regularization!
	\end{frame}
	
	\begin{frame}{Q51: Underflow Problem - Tutorial}
		\fontsize{6.5}{7.5}\selectfont
		\textbf{QUESTION: What is the underflow problem in Naïve Bayes?}
		
		\vspace{0.1cm}
		\textbf{OPTIONS:}
		\begin{enumerate}
			\item When probabilities become too large
			\item When multiplying many small probabilities results in numbers too small for computers to handle
			\item When documents are too long
			\item When vocabulary is too small
		\end{enumerate}
		
		\vspace{0.1cm}
		\textcolor{myblue}{\textbf{ANSWER: B}}
		
		\vspace{0.1cm}
		\textbf{DETAILED TUTORIAL:}
		\begin{itemize}
			\item \textbf{Underflow definition:}
			\begin{itemize}
				\item Multiplying many probabilities $< 1$
				\item Product becomes extremely small
				\item Below computer precision
			\end{itemize}
			\item \textbf{Solution:}
			\begin{itemize}
				\item Take logarithms
				\item $\log(P_1 \times ...) = \sum \log(P_i)$
			\end{itemize}
		\end{itemize}
		
		\textcolor{myred}{\textbf{EXAM TRAP:}} Underflow = \textcolor{myblue}{multiplying small probabilities} - use logarithms!
	\end{frame}
	
	\begin{frame}{Q52: Log Transformation - Tutorial}
		\fontsize{6.5}{7.5}\selectfont
		\textbf{QUESTION: Why is log transformation used in Naïve Bayes?}
		
		\vspace{0.1cm}
		\textbf{OPTIONS:}
		\begin{enumerate}
			\item To make probabilities larger
			\item To convert products to sums and avoid underflow
			\item To remove stop words
			\item To normalize the data
		\end{enumerate}
		
		\vspace{0.1cm}
		\textcolor{myblue}{\textbf{ANSWER: B}}
		
		\vspace{0.1cm}
		\textbf{DETAILED TUTORIAL:}
		\begin{itemize}
			\item \textbf{Log transformation benefits:}
			\begin{enumerate}
				\item \textbf{Products to sums:} $\log(P_1 \times P_2) = \log(P_1) + \log(P_2)$
				\item \textbf{Preserves ordering:} Class with highest probability = highest log probability
			\end{enumerate}
		\end{itemize}
		
		\textcolor{myred}{\textbf{EXAM TRAP:}} Log = \textcolor{myblue}{products to sums} - avoids underflow!
	\end{frame}
	
	\begin{frame}{Q53: Unknown Words Handling - Tutorial}
		\fontsize{6.5}{7.5}\selectfont
		\textbf{QUESTION: How are words in test data that were not in training data handled?}
		
		\vspace{0.1cm}
		\textbf{OPTIONS:}
		\begin{enumerate}
			\item They are added to the vocabulary
			\item They are ignored because they cannot help classify the document
			\item They cause the classification to fail
			\item They are treated as stop words
		\end{enumerate}
		
		\vspace{0.1cm}
		\textcolor{myblue}{\textbf{ANSWER: B}}
		
		\vspace{0.1cm}
		\textbf{DETAILED TUTORIAL:}
		\begin{itemize}
			\item \textbf{Unknown word handling:}
			\begin{itemize}
				\item Words not in training vocabulary
				\item No information about their distribution
				\item Must be removed from test document
			\end{itemize}
		\end{itemize}
		
		\textcolor{myred}{\textbf{EXAM TRAP:}} Unknown words = \textcolor{myblue}{ignored} - they provide no classification information!
	\end{frame}
	
	\begin{frame}{Q54: Naïve Bayes Assumptions - Tutorial}
		\fontsize{6.5}{7.5}\selectfont
		\textbf{QUESTION: What are the two key assumptions of Naïve Bayes for document classification?}
		
		\vspace{0.1cm}
		\textbf{OPTIONS:}
		\begin{enumerate}
			\item Words are dependent and have different weights
			\item Words are independent and have equal weighting
			\item Words are ordered and have different weights
			\item Words are random and have no weights
		\end{enumerate}
		
		\vspace{0.1cm}
		\textcolor{myblue}{\textbf{ANSWER: B}}
		
		\vspace{0.1cm}
		\textbf{DETAILED TUTORIAL:}
		\begin{itemize}
			\item \textbf{Assumption 1: Independence}
			\begin{itemize}
				\item Each word is independent of others
				\item $P(w_1, w_2, ..., w_n|c) = \prod P(w_i|c)$
			\end{itemize}
			\item \textbf{Assumption 2: Equal weighting}
			\begin{itemize}
				\item Each word has equal information value
			\end{itemize}
		\end{itemize}
		
		\textcolor{myred}{\textbf{EXAM TRAP:}} Naïve Bayes = \textcolor{myblue}{independence + equal weighting}!
	\end{frame}
	
	\begin{frame}{Q55: Generative vs Discriminative - Tutorial}
		\fontsize{6.5}{7.5}\selectfont
		\textbf{QUESTION: What makes Naïve Bayes a generative classifier?}
		
		\vspace{0.1cm}
		\textbf{OPTIONS:}
		\begin{enumerate}
			\item It can generate new documents
			\item It learns the joint distribution $P(d,c)$ and can create new classifications
			\item It only learns decision boundaries
			\item It doesn't use probabilities
		\end{enumerate}
		
		\vspace{0.1cm}
		\textcolor{myblue}{\textbf{ANSWER: B}}
		
		\vspace{0.1cm}
		\textbf{DETAILED TUTORIAL:}
		\begin{itemize}
			\item \textbf{Generative classifiers:}
			\begin{itemize}
				\item Model the joint distribution $P(d,c)$
				\item Can generate new examples
				\item Examples: Naïve Bayes
			\end{itemize}
		\end{itemize}
		
		\textcolor{myred}{\textbf{EXAM TRAP:}} Naïve Bayes = \textcolor{myblue}{generative} - learns joint distribution!
	\end{frame}
	
	\begin{frame}{Q56: Customer Service Example - Tutorial}
		\fontsize{6.5}{7.5}\selectfont
		\textbf{QUESTION: In the customer service example, what classification is made?}
		
		\vspace{0.1cm}
		\textbf{OPTIONS:}
		\begin{enumerate}
			\item "Good" because it has the highest prior
			\item "Bad" because it has the highest posterior probability
			\item "Indifferent" because it's the middle class
			\item Cannot be classified
		\end{enumerate}
		
		\vspace{0.1cm}
		\textcolor{myblue}{\textbf{ANSWER: B}}
		
		\vspace{0.1cm}
		\textbf{DETAILED TUTORIAL:}
		\begin{itemize}
			\item \textbf{Data:}
			\begin{itemize}
				\item 4 good, 2 bad, 2 indifferent (8 total)
				\item New document: words 1 and 2 present, words 3 and 4 absent
			\end{itemize}
			\item \textbf{Calculated probabilities:}
			\begin{itemize}
				\item $P(good|words) = 0$ (without smoothing)
				\item $P(bad|words) = 0.67$ (highest)
				\item $P(indifferent|words) = 0.33$
			\end{itemize}
		\end{itemize}
		
		\textcolor{myred}{\textbf{EXAM TRAP:}} Classification = \textcolor{myblue}{highest posterior probability} - not just prior!
	\end{frame}
	
	\begin{frame}{Q57: Smoothing Effect Example - Tutorial}
		\fontsize{6.5}{7.5}\selectfont
		\textbf{QUESTION: How did Laplace smoothing change the classification?}
		
		\vspace{0.1cm}
		\textbf{OPTIONS:}
		\begin{enumerate}
			\item It didn't change the classification
			\item It changed "Good" from 0 to 2.9\% probability
			\item It made "Bad" the lowest probability
			\item It made classification impossible
		\end{enumerate}
		
		\vspace{0.1cm}
		\textcolor{myblue}{\textbf{ANSWER: B}}
		
		\vspace{0.1cm}
		\textbf{DETAILED TUTORIAL:}
		\begin{itemize}
			\item \textbf{Before smoothing:}
			\begin{itemize}
				\item $P(good|words) = 0$
				\item $P(bad|words) = 0.67$
				\item $P(indifferent|words) = 0.33$
			\end{itemize}
			\item \textbf{After smoothing ($\lambda = 1$):}
			\begin{itemize}
				\item $P(good|words) = 0.029$ (increased from 0)
				\item $P(bad|words) = 0.583$ (decreased)
				\item $P(indifferent|words) = 0.388$ (increased slightly)
			\end{itemize}
		\end{itemize}
		
		\textcolor{myred}{\textbf{EXAM TRAP:}} Smoothing \textcolor{myblue}{eliminates zero probabilities} but may not change classification!
	\end{frame}
	
	\begin{frame}{Q58: Multi-class Classification - Tutorial}
		\fontsize{6.5}{7.5}\selectfont
		\textbf{QUESTION: How does Naïve Bayes handle more than two classes?}
		
		\vspace{0.1cm}
		\textbf{OPTIONS:}
		\begin{enumerate}
			\item It cannot handle more than two classes
			\item It computes posterior probability for each class and selects the highest
			\item It only uses binary classification
			\item It combines classes
		\end{enumerate}
		
		\vspace{0.1cm}
		\textcolor{myblue}{\textbf{ANSWER: B}}
		
		\vspace{0.1cm}
		\textbf{DETAILED TUTORIAL:}
		\begin{itemize}
			\item \textbf{Multi-class extension:}
			\begin{itemize}
				\item Natural extension of binary case
				\item Compute $P(c|d)$ for each class
				\item Select class with highest probability
			\end{itemize}
		\end{itemize}
		
		\textcolor{myred}{\textbf{EXAM TRAP:}} Multi-class = \textcolor{myblue}{choose class with highest posterior}!
	\end{frame}
	
	\begin{frame}{Q59: Naïve Bayes Advantages - Tutorial}
		\fontsize{6.5}{7.5}\selectfont
		\textbf{QUESTION: What are the main advantages of Naïve Bayes?}
		
		\vspace{0.1cm}
		\textbf{OPTIONS:}
		\begin{enumerate}
			\item It always has 100\% accuracy
			\item Simplicity, speed, and solid classification accuracy
			\item It works without any assumptions
			\item It requires no training data
		\end{enumerate}
		
		\vspace{0.1cm}
		\textcolor{myblue}{\textbf{ANSWER: B}}
		
		\vspace{0.1cm}
		\textbf{DETAILED TUTORIAL:}
		\begin{itemize}
			\item \textbf{Key advantages:}
			\begin{enumerate}
				\item \textbf{Simplicity:} Easy to understand and implement
				\item \textbf{Speed:} Fast training and prediction
				\item \textbf{Accuracy:} Solid classification despite assumptions
			\end{enumerate}
		\end{itemize}
		
		\textcolor{myred}{\textbf{EXAM TRAP:}} Naïve Bayes advantages = \textcolor{myblue}{simple, fast, and accurate}!
	\end{frame}
	
	\begin{frame}{Q60: Appendix Problem Classification - Tutorial}
		\fontsize{6.5}{7.5}\selectfont
		\textbf{QUESTION: In the appendix problem, what classification is made for the new instance?}
		
		\vspace{0.1cm}
		\textbf{OPTIONS:}
		\begin{enumerate}
			\item Buy (probability = 0.98)
			\item Sell (probability = 0.02)
			\item Cannot be classified
			\item Both equally likely
		\end{enumerate}
		
		\vspace{0.1cm}
		\textcolor{myblue}{\textbf{ANSWER: A}}
		
		\vspace{0.1cm}
		\textbf{DETAILED TUTORIAL:}
		\begin{itemize}
			\item \textbf{Data:}
			\begin{itemize}
				\item 5 buys, 5 sells
				\item Prior: $P(buy) = 0.5$, $P(sell) = 0.5$
				\item New document: words 1, 2, 3 present; words 4, 5 absent
			\end{itemize}
			\item \textbf{Calculations:}
			\begin{itemize}
				\item $P(buy) \times \prod P(w_i|buy) = 0.0768$
				\item $P(sell) \times \prod P(w_i|sell) = 0.00192$
			\end{itemize}
			\item \textbf{Posterior:} $P(buy|words) = 0.98$, $P(sell|words) = 0.02$
		\end{itemize}
		
		\textcolor{myred}{\textbf{EXAM TRAP:}} Work through the full calculation - \textcolor{myblue}{buy is much more likely}!
	\end{frame}
	
	% ============================================================
	% SECTION 5: WORD EMBEDDINGS & ADVANCED - QUESTIONS 61-70
	% ============================================================
	
	\begin{frame}{Q61: Word Embedding Definition - Tutorial}
		\fontsize{6.5}{7.5}\selectfont
		\textbf{QUESTION: What is word embedding in NLP?}
		
		\vspace{0.1cm}
		\textbf{OPTIONS:}
		\begin{enumerate}
			\item A method to remove words
			\item A technique to capture word meaning by representing words in a continuous vector space
			\item A method to count word frequencies
			\item A technique to stem words
		\end{enumerate}
		
		\vspace{0.1cm}
		\textcolor{myblue}{\textbf{ANSWER: B}}
		
		\vspace{0.1cm}
		\textbf{DETAILED TUTORIAL:}
		\begin{itemize}
			\item \textbf{Word embedding definition:}
			\begin{itemize}
				\item Represents words as dense vectors
				\item Captures semantic meaning
				\item Similar words have similar vectors
			\end{itemize}
		\end{itemize}
		
		\textcolor{myred}{\textbf{EXAM TRAP:}} Word embeddings = \textcolor{myblue}{capture meaning in dense vectors}!
	\end{frame}
	
	\begin{frame}{Q62: Word2Vec Definition - Tutorial}
		\fontsize{6.5}{7.5}\selectfont
		\textbf{QUESTION: What is Word2Vec?}
		
		\vspace{0.1cm}
		\textbf{OPTIONS:}
		\begin{enumerate}
			\item A word stemming algorithm
			\item A Google-developed algorithm for generating word embeddings using neural networks
			\item A dictionary approach to NLP
			\item A pre-processing technique
		\end{enumerate}
		
		\vspace{0.1cm}
		\textcolor{myblue}{\textbf{ANSWER: B}}
		
		\vspace{0.1cm}
		\textbf{DETAILED TUTORIAL:}
		\begin{itemize}
			\item \textbf{Word2Vec definition:}
			\begin{itemize}
				\item Developed by Google
				\item Neural network-based algorithm
				\item Generates word embeddings
			\end{itemize}
			\item \textbf{Two architectures:}
			\begin{enumerate}
				\item \textbf{CBOW:} Predict target from context
				\item \textbf{Skip-gram:} Predict context from target
			\end{enumerate}
		\end{itemize}
		
		\textcolor{myred}{\textbf{EXAM TRAP:}} Word2Vec = \textcolor{myblue}{Google algorithm for word embeddings}!
	\end{frame}
	
	\begin{frame}{Q63: Context in Word Embeddings - Tutorial}
		\fontsize{6.5}{7.5}\selectfont
		\textbf{QUESTION: How do word embeddings capture context?}
		
		\vspace{0.1cm}
		\textbf{OPTIONS:}
		\begin{enumerate}
			\item By ignoring context
			\item By paying attention to positions of words and their surrounding words
			\item By counting all words equally
			\item By removing all punctuation
		\end{enumerate}
		
		\vspace{0.1cm}
		\textcolor{myblue}{\textbf{ANSWER: B}}
		
		\vspace{0.1cm}
		\textbf{DETAILED TUTORIAL:}
		\begin{itemize}
			\item \textbf{Context capture methods:}
			\begin{enumerate}
				\item \textbf{Position information:} Track which words appear together
				\item \textbf{Surrounding words:} Consider words before and after
			\end{enumerate}
		\end{itemize}
		
		\textcolor{myred}{\textbf{EXAM TRAP:}} Word embeddings capture context by \textcolor{myblue}{analyzing surrounding words}!
	\end{frame}
	
	\begin{frame}{Q64: Skip-gram vs CBOW - Tutorial}
		\fontsize{6.5}{7.5}\selectfont
		\textbf{QUESTION: What is the difference between Skip-gram and CBOW?}
		
		\vspace{0.1cm}
		\textbf{OPTIONS:}
		\begin{enumerate}
			\item They are the same algorithm
			\item Skip-gram predicts context from target; CBOW predicts target from context
			\item Skip-gram is for English only
			\item CBOW is for English only
		\end{enumerate}
		
		\vspace{0.1cm}
		\textcolor{myblue}{\textbf{ANSWER: B}}
		
		\vspace{0.1cm}
		\textbf{DETAILED TUTORIAL:}
		\begin{itemize}
			\item \textbf{CBOW:} Context words $\to$ target word (faster, better for frequent words)
			\item \textbf{Skip-gram:} Target word $\to$ context words (better for rare words, more expensive)
		\end{itemize}
		
		\textcolor{myred}{\textbf{EXAM TRAP:}} Skip-gram = \textcolor{myblue}{target $\to$ context}; CBOW = \textcolor{myblue}{context $\to$ target}!
	\end{frame}
	
	\begin{frame}{Q65: One-Hot Encoding - Tutorial}
		\fontsize{6.5}{7.5}\selectfont
		\textbf{QUESTION: What is one-hot encoding in NLP?}
		
		\vspace{0.1cm}
		\textbf{OPTIONS:}
		\begin{enumerate}
			\item A technique to encode temperature
			\item Representing words as binary vectors with a 1 in the position of that word and 0 elsewhere
			\item A method to count word frequencies
			\item A technique to remove stop words
		\end{enumerate}
		
		\vspace{0.1cm}
		\textcolor{myblue}{\textbf{ANSWER: B}}
		
		\vspace{0.1cm}
		\textbf{DETAILED TUTORIAL:}
		\begin{itemize}
			\item \textbf{One-hot encoding:}
			\begin{itemize}
				\item Vector length = vocabulary size
				\item One position = 1, all others = 0
				\item Very sparse, no semantic meaning
			\end{itemize}
		\end{itemize}
		
		\textcolor{myred}{\textbf{EXAM TRAP:}} One-hot = \textcolor{myblue}{binary vector with one 1} - no semantic meaning!
	\end{frame}
	
	\begin{frame}{Q66: Dimensionality Reduction in NLP - Tutorial}
		\fontsize{6.5}{7.5}\selectfont
		\textbf{QUESTION: Why is dimensionality reduction important in NLP?}
		
		\vspace{0.1cm}
		\textbf{OPTIONS:}
		\begin{enumerate}
			\item To make documents longer
			\item To reduce sparse, high-dimensional vectors to dense, lower-dimensional representations
			\item To add more features
			\item To increase computation time
		\end{enumerate}
		
		\vspace{0.1cm}
		\textcolor{myblue}{\textbf{ANSWER: B}}
		
		\vspace{0.1cm}
		\textbf{DETAILED TUTORIAL:}
		\begin{itemize}
			\item \textbf{The problem:}
			\begin{itemize}
				\item Vocabulary can have 10,000+ words
				\item Very sparse, slow, and inefficient
			\end{itemize}
			\item \textbf{Reduction methods:} PCA, SVD, Word2Vec
		\end{itemize}
		
		\textcolor{myred}{\textbf{EXAM TRAP:}} Dimensionality reduction = \textcolor{myblue}{sparse $\to$ dense} representations!
	\end{frame}
	
	\begin{frame}{Q67: BoW vs Word Embeddings - Tutorial}
		\fontsize{6.5}{7.5}\selectfont
		\textbf{QUESTION: What is a key advantage of word embeddings over BoW?}
		
		\vspace{0.1cm}
		\textbf{OPTIONS:}
		\begin{enumerate}
			\item Embeddings are sparser
			\item Embeddings capture semantic meaning and relationships between words
			\item Embeddings are faster to compute
			\item Embeddings don't require pre-processing
		\end{enumerate}
		
		\vspace{0.1cm}
		\textcolor{myblue}{\textbf{ANSWER: B}}
		
		\vspace{0.1cm}
		\textbf{DETAILED TUTORIAL:}
		\begin{itemize}
			\item \textbf{Embedding advantages:}
			\begin{enumerate}
				\item \textbf{Semantic meaning:} Similar words close in space
				\item \textbf{Dense vectors:} Lower dimensional, more efficient
				\item \textbf{Context awareness:} Considers surrounding words
			\end{enumerate}
		\end{itemize}
		
		\textcolor{myred}{\textbf{EXAM TRAP:}} Embeddings = \textcolor{myblue}{capture semantic meaning} - BoW doesn't!
	\end{frame}
	
	\begin{frame}{Q68: Skip-gram Details - Tutorial}
		\fontsize{6.5}{7.5}\selectfont
		\textbf{QUESTION: What is Skip-gram used for in NLP?}
		
		\vspace{0.1cm}
		\textbf{OPTIONS:}
		\begin{enumerate}
			\item To remove stop words
			\item To predict surrounding words given a target word
			\item To count word frequencies
			\item To stem words
		\end{enumerate}
		
		\vspace{0.1cm}
		\textcolor{myblue}{\textbf{ANSWER: B}}
		
		\vspace{0.1cm}
		\textbf{DETAILED TUTORIAL:}
		\begin{itemize}
			\item \textbf{Skip-gram architecture:}
			\begin{itemize}
				\item Input: target word
				\item Output: context words
				\item Uses neural network
			\end{itemize}
		\end{itemize}
		
		\textcolor{myred}{\textbf{EXAM TRAP:}} Skip-gram = \textcolor{myblue}{target word $\to$ predict context words}!
	\end{frame}
	
	\begin{frame}{Q69: Word2Vec Application - Tutorial}
		\fontsize{6.5}{7.5}\selectfont
		\textbf{QUESTION: What is a famous semantic relationship captured by Word2Vec?}
		
		\vspace{0.1cm}
		\textbf{OPTIONS:}
		\begin{enumerate}
			\item "king - man + woman equals queen"
			\item "king + man + woman equals queen"
			\item "king - queen equals man - woman"
			\item All words are equally related
		\end{enumerate}
		
		\vspace{0.1cm}
		\textcolor{myblue}{\textbf{ANSWER: A}}
		
		\vspace{0.1cm}
		\textbf{DETAILED TUTORIAL:}
		\begin{itemize}
			\item \textbf{Famous analogy:}
			\begin{itemize}
				\item $v_{king} - v_{man} + v_{woman} \approx v_{queen}$
				\item Captures gender relationship
				\item Preserves arithmetic operations
			\end{itemize}
		\end{itemize}
		
		\textcolor{myred}{\textbf{EXAM TRAP:}} Word2Vec captures \textcolor{myblue}{semantic relationships} via vector arithmetic!
	\end{frame}
	
	\begin{frame}{Q70: NLP Evaluation Overview - Tutorial}
		\fontsize{6.5}{7.5}\selectfont
		\textbf{QUESTION: How are NLP models evaluated?}
		
		\vspace{0.1cm}
		\textbf{OPTIONS:}
		\begin{enumerate}
			\item Only by accuracy
			\item Using confusion matrices, accuracy, precision, recall, F1 when labels exist
			\item Only by speed
			\item Only by data requirements
		\end{enumerate}
		
		\vspace{0.1cm}
		\textcolor{myblue}{\textbf{ANSWER: B}}
		
		\vspace{0.1cm}
		\textbf{DETAILED TUTORIAL:}
		\begin{itemize}
			\item \textbf{Labeled data evaluation:}
			\begin{enumerate}
				\item \textbf{Confusion Matrix:} TP, TN, FP, FN
				\item \textbf{Accuracy:} Overall correctness
				\item \textbf{Precision:} Of predicted positives, how many were right?
				\item \textbf{Recall:} Of actual positives, how many did we catch?
				\item \textbf{F1:} Harmonic mean of precision and recall
			\end{enumerate}
		\end{itemize}
		
		\textcolor{myred}{\textbf{EXAM TRAP:}} Evaluation = \textcolor{myblue}{confusion matrix metrics} when labels exist!
	\end{frame}
	
	% ============================================================
	% SECTION 6: EVALUATION & PRACTICAL - QUESTIONS 71-80
	% ============================================================
	
	\begin{frame}{Q71: NLP Model Evaluation Metrics - Tutorial}
		\fontsize{6.5}{7.5}\selectfont
		\textbf{QUESTION: What metrics are used to evaluate NLP classification models?}
		
		\vspace{0.1cm}
		\textbf{OPTIONS:}
		\begin{enumerate}
			\item Only accuracy
			\item Accuracy, precision, recall, and F1 from confusion matrix
			\item Only precision
			\item Only recall
		\end{enumerate}
		
		\vspace{0.1cm}
		\textcolor{myblue}{\textbf{ANSWER: B}}
		
		\vspace{0.1cm}
		\textbf{DETAILED TUTORIAL:}
		\begin{itemize}
			\item \textbf{Standard evaluation metrics:}
			\begin{enumerate}
				\item \textbf{Accuracy:} $(TP + TN) / (TP + TN + FP + FN)$
				\item \textbf{Precision:} $TP / (TP + FP)$
				\item \textbf{Recall:} $TP / (TP + FN)$
				\item \textbf{F1:} Harmonic mean of precision and recall
			\end{enumerate}
		\end{itemize}
		
		\textcolor{myred}{\textbf{EXAM TRAP:}} Use \textcolor{myblue}{multiple metrics} - not just accuracy!
	\end{frame}
	
	\begin{frame}{Q72: Labeled vs Unlabeled Data - Tutorial}
		\fontsize{6.5}{7.5}\selectfont
		\textbf{QUESTION: How is NLP evaluation different for unlabeled data?}
		
		\vspace{0.1cm}
		\textbf{OPTIONS:}
		\begin{enumerate}
			\item No evaluation is possible
			\item Use unsupervised evaluation methods or qualitative assessment
			\item Use the same metrics as labeled data
			\item Evaluation is not needed
		\end{enumerate}
		
		\vspace{0.1cm}
		\textcolor{myblue}{\textbf{ANSWER: B}}
		
		\vspace{0.1cm}
		\textbf{DETAILED TUTORIAL:}
		\begin{itemize}
			\item \textbf{Unlabeled data evaluation:}
			\begin{itemize}
				\item Subject matter expert review
				\item Unsupervised metrics (cohesion, separation)
				\item Indirect validation
			\end{itemize}
		\end{itemize}
		
		\textcolor{myred}{\textbf{EXAM TRAP:}} Unlabeled data requires \textcolor{myblue}{alternative evaluation methods}!
	\end{frame}
	
	\begin{frame}{Q73: Speed and Data Requirements - Tutorial}
		\fontsize{6.5}{7.5}\selectfont
		\textbf{QUESTION: Why should speed and data requirements be considered in NLP evaluation?}
		
		\vspace{0.1cm}
		\textbf{OPTIONS:}
		\begin{enumerate}
			\item They are not important
			\item Textual databases can be vast and sparse vectors are hard to handle efficiently
			\item Speed is never a concern
			\item All NLP models are equally fast
		\end{enumerate}
		
		\vspace{0.1cm}
		\textcolor{myblue}{\textbf{ANSWER: B}}
		
		\vspace{0.1cm}
		\textbf{DETAILED TUTORIAL:}
		\begin{itemize}
			\item \textbf{Practical considerations:}
			\begin{enumerate}
				\item \textbf{Data volume:} NLP datasets can be enormous
				\item \textbf{Sparsity:} Sparse vectors are inefficient
				\item \textbf{Model complexity:} Neural networks are slow to train
			\end{enumerate}
		\end{itemize}
		
		\textcolor{myred}{\textbf{EXAM TRAP:}} NLP can be \textcolor{myblue}{computationally expensive} - speed and data matter!
	\end{frame}
	
	\begin{frame}{Q74: BoW Limitations - Tutorial}
		\fontsize{6.5}{7.5}\selectfont
		\textbf{QUESTION: What is a major limitation of simple BoW methods?}
		
		\vspace{0.1cm}
		\textbf{OPTIONS:}
		\begin{enumerate}
			\item They are too fast
			\item They are not useful for determining context or intended meaning
			\item They are too accurate
			\item They require too much data
		\end{enumerate}
		
		\vspace{0.1cm}
		\textcolor{myblue}{\textbf{ANSWER: B}}
		
		\vspace{0.1cm}
		\textbf{DETAILED TUTORIAL:}
		\begin{itemize}
			\item \textbf{BoW limitations:}
			\begin{enumerate}
				\item \textbf{Loss of context:} Word order ignored
				\item \textbf{No semantic meaning:} No concept of similarity
				\item \textbf{Negation issues:} "not good" vs "good"
			\end{enumerate}
		\end{itemize}
		
		\textcolor{myred}{\textbf{EXAM TRAP:}} BoW fails to capture \textcolor{myblue}{context and meaning}!
	\end{frame}
	
	\begin{frame}{Q75: Advanced NLP Techniques - Tutorial}
		\fontsize{6.5}{7.5}\selectfont
		\textbf{QUESTION: What advanced techniques help determine word context?}
		
		\vspace{0.1cm}
		\textbf{OPTIONS:}
		\begin{enumerate}
			\item Only BoW
			\item n-grams, Skip-grams, and word embeddings
			\item Only stop word removal
			\item Only stemming
		\end{enumerate}
		
		\vspace{0.1cm}
		\textcolor{myblue}{\textbf{ANSWER: B}}
		
		\vspace{0.1cm}
		\textbf{DETAILED TUTORIAL:}
		\begin{itemize}
			\item \textbf{Context-aware techniques:}
			\begin{enumerate}
				\item \textbf{n-grams:} n contiguous words as one unit
				\item \textbf{Skip-grams:} Predict context from target
				\item \textbf{Word embeddings:} Dense vector representations
			\end{enumerate}
		\end{itemize}
		
		\textcolor{myred}{\textbf{EXAM TRAP:}} \textcolor{myblue}{n-grams, Skip-grams, and embeddings} capture context!
	\end{frame}
	
	\begin{frame}{Q76: TF-IDF vs Word Counts - Tutorial}
		\fontsize{6.5}{7.5}\selectfont
		\textbf{QUESTION: What is the advantage of TF-IDF over simple word counts?}
		
		\vspace{0.1cm}
		\textbf{OPTIONS:}
		\begin{enumerate}
			\item TF-IDF is faster to compute
			\item TF-IDF assigns higher weights to rare but important words
			\item TF-IDF counts all words equally
			\item TF-IDF doesn't require a corpus
		\end{enumerate}
		
		\vspace{0.1cm}
		\textcolor{myblue}{\textbf{ANSWER: B}}
		
		\vspace{0.1cm}
		\textbf{DETAILED TUTORIAL:}
		\begin{itemize}
			\item \textbf{TF-IDF advantages:}
			\begin{enumerate}
				\item \textbf{Weights rare words higher:} More discriminatory power
				\item \textbf{Weights common words lower:} Reduced noise
			\end{enumerate}
		\end{itemize}
		
		\textcolor{myred}{\textbf{EXAM TRAP:}} TF-IDF = \textcolor{myblue}{higher weight for rare words}!
	\end{frame}
	
	\begin{frame}{Q77: Document Similarity Measurement - Tutorial}
		\fontsize{6.5}{7.5}\selectfont
		\textbf{QUESTION: How is document similarity measured in NLP?}
		
		\vspace{0.1cm}
		\textbf{OPTIONS:}
		\begin{enumerate}
			\item By comparing document titles
			\item By comparing document vectors using similarity metrics (cosine similarity)
			\item By counting words in each document
			\item By comparing document lengths
		\end{enumerate}
		
		\vspace{0.1cm}
		\textcolor{myblue}{\textbf{ANSWER: B}}
		
		\vspace{0.1cm}
		\textbf{DETAILED TUTORIAL:}
		\begin{itemize}
			\item \textbf{Similarity measurement:}
			\begin{enumerate}
				\item \textbf{Document vectors:} Using BoW, TF-IDF, or embeddings
				\item \textbf{Similarity metrics:} Cosine similarity, Jaccard, Euclidean
			\end{enumerate}
		\end{itemize}
		
		\textcolor{myred}{\textbf{EXAM TRAP:}} Document similarity = \textcolor{myblue}{comparing document vectors} - often cosine similarity!
	\end{frame}
	
	\begin{frame}{Q78: Cosine Similarity - Tutorial}
		\fontsize{6.5}{7.5}\selectfont
		\textbf{QUESTION: What is cosine similarity used for in NLP?}
		
		\vspace{0.1cm}
		\textbf{OPTIONS:}
		\begin{enumerate}
			\item To count word frequencies
			\item To measure similarity between two document vectors
			\item To remove stop words
			\item To stem words
		\end{enumerate}
		
		\vspace{0.1cm}
		\textcolor{myblue}{\textbf{ANSWER: B}}
		
		\vspace{0.1cm}
		\textbf{DETAILED TUTORIAL:}
		\begin{itemize}
			\item \textbf{Cosine similarity formula:}
			\begin{itemize}
				\item $\cos(\theta) = \frac{v_1 \cdot v_2}{||v_1|| \times ||v_2||}$
				\item Range: -1 to 1
				\item cos equals 1 = very similar
			\end{itemize}
		\end{itemize}
		
		\textcolor{myred}{\textbf{EXAM TRAP:}} Cosine similarity = \textcolor{myblue}{measures document vector similarity}!
	\end{frame}
	
	\begin{frame}{Q79: Stemming vs Lemmatization Tradeoff - Tutorial}
		\fontsize{6.5}{7.5}\selectfont
		\textbf{QUESTION: What is the tradeoff between stemming and lemmatization?}
		
		\vspace{0.1cm}
		\textbf{OPTIONS:}
		\begin{enumerate}
			\item No tradeoff - they are the same
			\item Stemming is faster but less accurate; lemmatization is slower but more accurate
			\item Stemming is more accurate
			\item Lemmatization is faster
		\end{enumerate}
		
		\vspace{0.1cm}
		\textcolor{myblue}{\textbf{ANSWER: B}}
		
		\vspace{0.1cm}
		\textbf{DETAILED TUTORIAL:}
		\begin{itemize}
			\item \textbf{Stemming:} Fast, less accurate, may create non-words
			\item \textbf{Lemmatization:} Slow, more accurate, always produces real words
		\end{itemize}
		
		\textcolor{myred}{\textbf{EXAM TRAP:}} Stemming = \textcolor{myblue}{fast but less accurate}; Lemmatization = \textcolor{myblue}{slow but more accurate}!
	\end{frame}
	
	\begin{frame}{Q80: NLP Challenges Summary - Tutorial}
		\fontsize{6.5}{7.5}\selectfont
		\textbf{QUESTION: What are the main challenges in NLP?}
		
		\vspace{0.1cm}
		\textbf{OPTIONS:}
		\begin{enumerate}
			\item Only computational speed
			\item Ambiguity, context dependence, noise, and sarcasm/detecting tone
			\item Only data size
			\item Only language barriers
		\end{enumerate}
		
		\vspace{0.1cm}
		\textcolor{myblue}{\textbf{ANSWER: B}}
		
		\vspace{0.1cm}
		\textbf{DETAILED TUTORIAL:}
		\begin{itemize}
			\item \textbf{Key challenges:}
			\begin{enumerate}
				\item \textbf{Ambiguity:} Same words have different meanings
				\item \textbf{Context dependence:} Meaning depends on surrounding words
				\item \textbf{Noise:} Spelling errors, abbreviations, colloquialisms
				\item \textbf{Tone detection:} Sarcasm, sentiment nuances
			\end{enumerate}
		\end{itemize}
		
		\textcolor{myred}{\textbf{EXAM TRAP:}} NLP challenges include \textcolor{myblue}{ambiguity, context, noise, and tone}!
	\end{frame}
	
	% ============================================================
	% SECTION 7: CHAPTER QUESTIONS - QUESTIONS 81-100
	% ============================================================
	
	\begin{frame}{Q81: Sentiment Analysis Steps - Tutorial}
		\fontsize{6.5}{7.5}\selectfont
		\textbf{QUESTION: What are the basic steps for sentiment analysis using a dictionary approach?}
		
		\vspace{0.1cm}
		\textbf{OPTIONS:}
		\begin{enumerate}
			\item Only count positive words
			\item Establish dictionary, pre-process text, replace with stems/lemmas, calculate proportions
			\item Only count negative words
			\item Use machine learning only
		\end{enumerate}
		
		\vspace{0.1cm}
		\textcolor{myblue}{\textbf{ANSWER: B}}
		
		\vspace{0.1cm}
		\textbf{DETAILED TUTORIAL:}
		\begin{itemize}
			\item \textbf{Sentiment analysis steps:}
			\begin{enumerate}
				\item \textbf{Establish dictionary:} Positive, negative, neutral lists
				\item \textbf{Pre-process text:} Remove punctuation, stop words, lowercase
				\item \textbf{Replace with stems/lemmas:} Standardize word forms
				\item \textbf{Calculate proportions:} Positive, negative, neutral
			\end{enumerate}
		\end{itemize}
		
		\textcolor{myred}{\textbf{EXAM TRAP:}} Steps = \textcolor{myblue}{dictionary, pre-process, stem, calculate proportions}!
	\end{frame}
	
	\begin{frame}{Q82: BoW Model Explanation - Tutorial}
		\fontsize{6.5}{7.5}\selectfont
		\textbf{QUESTION: Explain how a bag of words model is used in NLP.}
		
		\vspace{0.1cm}
		\textbf{OPTIONS:}
		\begin{enumerate}
			\item It preserves word order
			\item It treats each term identically, counts frequencies, stores in vectors for analysis
			\item It removes all words
			\item It only works for numbers
		\end{enumerate}
		
		\vspace{0.1cm}
		\textcolor{myblue}{\textbf{ANSWER: B}}
		
		\vspace{0.1cm}
		\textbf{DETAILED TUTORIAL:}
		\begin{itemize}
			\item \textbf{BoW process:}
			\begin{enumerate}
				\item Treat terms identically (ignore order)
				\item Count frequencies of each word
				\item Store in vectors
			\end{enumerate}
		\end{itemize}
		
		\textcolor{myred}{\textbf{EXAM TRAP:}} BoW = \textcolor{myblue}{treats each word identically, counts, stores in vectors}!
	\end{frame}
	
	\begin{frame}{Q83: Naïve Bayes Assumptions - Tutorial}
		\fontsize{6.5}{7.5}\selectfont
		\textbf{QUESTION: What are the two assumptions of Naïve Bayes for document grouping?}
		
		\vspace{0.1cm}
		\textbf{OPTIONS:}
		\begin{enumerate}
			\item Words are ordered and have different weights
			\item Each word is independent of others and each word has equal weighting
			\item Words are dependent and have different weights
			\item Words are random
		\end{enumerate}
		
		\vspace{0.1cm}
		\textcolor{myblue}{\textbf{ANSWER: B}}
		
		\vspace{0.1cm}
		\textbf{DETAILED TUTORIAL:}
		\begin{itemize}
			\item \textbf{Assumption 1: Independence}
			\begin{itemize}
				\item $P(w_1, w_2, ..., w_n|c) = P(w_1|c) \times P(w_2|c) \times ...$
			\end{itemize}
			\item \textbf{Assumption 2: Equal weighting}
			\begin{itemize}
				\item Each word has equal information value
			\end{itemize}
		\end{itemize}
		
		\textcolor{myred}{\textbf{EXAM TRAP:}} Naïve Bayes = \textcolor{myblue}{independence + equal weighting} - two key assumptions!
	\end{frame}
	
	\begin{frame}{Q84: Negation Problem and Solution - Tutorial}
		\fontsize{6.5}{7.5}\selectfont
		\textbf{QUESTION: Why do negation words cause problems and how can they be handled?}
		
		\vspace{0.1cm}
		\textbf{OPTIONS:}
		\begin{enumerate}
			\item They don't cause problems
			\item Negations reverse meaning; use n-grams to pair negation with following words
			\item Remove all negation words
			\item Treat negation as positive
		\end{enumerate}
		
		\vspace{0.1cm}
		\textcolor{myblue}{\textbf{ANSWER: B}}
		
		\vspace{0.1cm}
		\textbf{DETAILED TUTORIAL:}
		\begin{itemize}
			\item \textbf{The negation problem:}
			\begin{itemize}
				\item "not adequate" vs "adequate"
				\item BoW sees "adequate" in both
				\item Meaning is opposite
			\end{itemize}
			\item \textbf{n-gram solution:}
			\begin{itemize}
				\item Treat "not adequate" as one unit
				\item Bi-grams capture negation
			\end{itemize}
		\end{itemize}
		
		\textcolor{myred}{\textbf{EXAM TRAP:}} Negation problem = \textcolor{myblue}{solved with n-grams} pairing negation + following word!
	\end{frame}
	
	\begin{frame}{Q85: Zero TF-IDF Explanation - Tutorial}
		\fontsize{6.5}{7.5}\selectfont
		\textbf{QUESTION: Why do some words have zero TF-IDF values?}
		
		\vspace{0.1cm}
		\textbf{OPTIONS:}
		\begin{enumerate}
			\item Because they are misspelled
			\item Because they appear in every document and have no discriminatory power
			\item Because they are too long
			\item Because they are not in the vocabulary
		\end{enumerate}
		
		\vspace{0.1cm}
		\textcolor{myblue}{\textbf{ANSWER: B}}
		
		\vspace{0.1cm}
		\textbf{DETAILED TUTORIAL:}
		\begin{itemize}
			\item \textbf{TF-IDF = 0 condition:}
			\begin{itemize}
				\item Word appears in all documents
				\item IDF = ln(D/D) = ln(1) = 0
				\item TF-IDF = TF $\times$ 0 = 0
			\end{itemize}
		\end{itemize}
		
		\textcolor{myred}{\textbf{EXAM TRAP:}} Zero TF-IDF = \textcolor{myblue}{word appears in all documents} - no discriminatory power!
	\end{frame}
	
	\begin{frame}{Q86: TF-IDF Variation by Document - Tutorial}
		\fontsize{6.5}{7.5}\selectfont
		\textbf{QUESTION: Why does the same word have different TF-IDF scores in different documents?}
		
		\vspace{0.1cm}
		\textbf{OPTIONS:}
		\begin{enumerate}
			\item Because IDF varies by document
			\item Because TF varies by document (document length and frequency differences)
			\item Because vocabulary changes
			\item Because documents have different authors
		\end{enumerate}
		
		\vspace{0.1cm}
		\textcolor{myblue}{\textbf{ANSWER: B}}
		
		\vspace{0.1cm}
		\textbf{DETAILED TUTORIAL:}
		\begin{itemize}
			\item \textbf{TF variation:}
			\begin{itemize}
				\item $TF = w_{i,j} / |v_j|$
				\item Different lengths $\Rightarrow$ different TFs
			\end{itemize}
			\item \textbf{IDF is constant:} Depends only on corpus
		\end{itemize}
		
		\textcolor{myred}{\textbf{EXAM TRAP:}} TF varies by document $\Rightarrow$ \textcolor{myblue}{TF-IDF varies by document}!
	\end{frame}
	
	\begin{frame}{Q87: Appendix Problem Unknown Words - Tutorial}
		\fontsize{6.5}{7.5}\selectfont
		\textbf{QUESTION: If new words appear in test documents not in the training sample, how should they be treated?}
		
		\vspace{0.1cm}
		\textbf{OPTIONS:}
		\begin{enumerate}
			\item Add them to the vocabulary
			\item Ignore them because they cannot help with classification
			\item Remove the entire document
			\item Use them as stop words
		\end{enumerate}
		
		\vspace{0.1cm}
		\textcolor{myblue}{\textbf{ANSWER: B}}
		
		\vspace{0.1cm}
		\textbf{DETAILED TUTORIAL:}
		\begin{itemize}
			\item \textbf{Unknown word handling:}
			\begin{itemize}
				\item Words not in training vocabulary
				\item No information about their distribution
				\item Cannot calculate $P(w|c)$
				\item Must be ignored
			\end{itemize}
		\end{itemize}
		
		\textcolor{myred}{\textbf{EXAM TRAP:}} Unknown words = \textcolor{myblue}{ignored} - no information about them!
	\end{frame}
	
	\begin{frame}{Q88: Smoothing Explanation - Tutorial}
		\fontsize{6.5}{7.5}\selectfont
		\textbf{QUESTION: Why is smoothing required in Naïve Bayes applications and how does it work?}
		
		\vspace{0.1cm}
		\textbf{OPTIONS:}
		\begin{enumerate}
			\item It's not required
			\item To avoid zero probability problems by adding a positive integer to all counts
			\item To remove stop words
			\item To stem words
		\end{enumerate}
		
		\vspace{0.1cm}
		\textcolor{myblue}{\textbf{ANSWER: B}}
		
		\vspace{0.1cm}
		\textbf{DETAILED TUTORIAL:}
		\begin{itemize}
			\item \textbf{Why smoothing is needed:}
			\begin{itemize}
				\item Zero probability problem
				\item Word not in training class $\to$ P = 0
				\item Product becomes zero
			\end{itemize}
			\item \textbf{How smoothing works:}
			\begin{itemize}
				\item Add $\lambda$ to all counts
				\item $\lambda = 1$: Laplace smoothing
			\end{itemize}
		\end{itemize}
		
		\textcolor{myred}{\textbf{EXAM TRAP:}} Smoothing = \textcolor{myblue}{add to counts to avoid zero probabilities}!
	\end{frame}
	
	\begin{frame}{Q89: Laplace Smoothing Impact - Tutorial}
		\fontsize{6.5}{7.5}\selectfont
		\textbf{QUESTION: What happens to the probabilities under Laplace smoothing?}
		
		\vspace{0.1cm}
		\textbf{OPTIONS:}
		\begin{enumerate}
			\item They become more extreme
			\item They are pushed toward a uniform distribution
			\item They become negative
			\item They don't change
		\end{enumerate}
		
		\vspace{0.1cm}
		\textcolor{myblue}{\textbf{ANSWER: B}}
		
		\vspace{0.1cm}
		\textbf{DETAILED TUTORIAL:}
		\begin{itemize}
			\item \textbf{Effect of Laplace smoothing:}
			\begin{itemize}
				\item Adds counts to all classes
				\item All probabilities become positive
				\item Shrinks toward uniform distribution
				\item Similar to regularization
			\end{itemize}
		\end{itemize}
		
		\textcolor{myred}{\textbf{EXAM TRAP:}} Laplace smoothing = \textcolor{myblue}{shrinks toward uniform distribution}!
	\end{frame}
	
	\begin{frame}{Q90: Appendix Problem Classification - Tutorial}
		\fontsize{6.5}{7.5}\selectfont
		\textbf{QUESTION: In the appendix problem, what is the most likely recommendation?}
		
		\vspace{0.1cm}
		\textbf{OPTIONS:}
		\begin{enumerate}
			\item Buy (probability = 0.98)
			\item Sell (probability = 0.02)
			\item Hold
			\item Cannot be determined
		\end{enumerate}
		
		\vspace{0.1cm}
		\textcolor{myblue}{\textbf{ANSWER: A}}
		
		\vspace{0.1cm}
		\textbf{DETAILED TUTORIAL:}
		\begin{itemize}
			\item \textbf{Given data:}
			\begin{itemize}
				\item 5 buys, 5 sells (prior = 0.5 each)
				\item New document: words 1-3 present, words 4-5 absent
			\end{itemize}
			\item \textbf{Calculations:}
			\begin{itemize}
				\item Buy: 0.8 $\times$ 1.0 $\times$ 0.4 $\times$ 0.8 $\times$ 0.6 $\times$ 0.5 = 0.0768
				\item Sell: 0.2 $\times$ 0.4 $\times$ 0.6 $\times$ 0.4 $\times$ 0.2 $\times$ 0.5 = 0.00192
			\end{itemize}
			\item \textbf{Result:} Buy (98\% probability)
		\end{itemize}
		
		\textcolor{myred}{\textbf{EXAM TRAP:}} Work through the full calculation - \textcolor{myblue}{buy is much more likely}!
	\end{frame}
	
	\begin{frame}{Q91: NLP Pre-processing Steps - Tutorial}
		\fontsize{6.5}{7.5}\selectfont
		\textbf{QUESTION: What are the main pre-processing steps in NLP?}
		
		\vspace{0.1cm}
		\textbf{OPTIONS:}
		\begin{enumerate}
			\item Only tokenization
			\item Tokenization, stop word removal, stemming/lemmatization, and feature extraction
			\item Only stemming
			\item Only stop word removal
		\end{enumerate}
		
		\vspace{0.1cm}
		\textcolor{myblue}{\textbf{ANSWER: B}}
		
		\vspace{0.1cm}
		\textbf{DETAILED TUTORIAL:}
		\begin{itemize}
			\item \textbf{Pre-processing steps:}
			\begin{enumerate}
				\item \textbf{Tokenization:} Split into sentences and words
				\item \textbf{Stop word removal:} Remove common words with no value
				\item \textbf{Stemming/Lemmatization:} Reduce to root form
				\item \textbf{Feature extraction:} Convert to numeric vectors
			\end{enumerate}
		\end{itemize}
		
		\textcolor{myred}{\textbf{EXAM TRAP:}} Pre-processing = \textcolor{myblue}{tokenization, stop words, stemming/lemmatization, feature extraction}!
	\end{frame}
	
	\begin{frame}{Q92: Tokenization Detail - Tutorial}
		\fontsize{6.5}{7.5}\selectfont
		\textbf{QUESTION: What is the difference between sentence tokenization and word tokenization?}
		
		\vspace{0.1cm}
		\textbf{OPTIONS:}
		\begin{enumerate}
			\item They are the same thing
			\item Sentence tokenization splits at punctuation; word tokenization splits into words
			\item Sentence tokenization removes words
			\item Word tokenization removes sentences
		\end{enumerate}
		
		\vspace{0.1cm}
		\textcolor{myblue}{\textbf{ANSWER: B}}
		
		\vspace{0.1cm}
		\textbf{DETAILED TUTORIAL:}
		\begin{itemize}
			\item \textbf{Sentence tokenization:} Splits at periods, question marks, exclamation marks
			\item \textbf{Word tokenization:} Splits sentences into words, removes punctuation
		\end{itemize}
		
		\textcolor{myred}{\textbf{EXAM TRAP:}} Sentence = \textcolor{myblue}{splitting at punctuation}; Word = \textcolor{myblue}{splitting into words}!
	\end{frame}
	
	\begin{frame}{Q93: Stop Words Importance - Tutorial}
		\fontsize{6.5}{7.5}\selectfont
		\textbf{QUESTION: Why do stop words vary across applications?}
		
		\vspace{0.1cm}
		\textbf{OPTIONS:}
		\begin{enumerate}
			\item All stop words are the same
			\item Words with no value in one context may be useful in another
			\item Stop words are always removed
			\item Stop words never matter
		\end{enumerate}
		
		\vspace{0.1cm}
		\textcolor{myblue}{\textbf{ANSWER: B}}
		
		\vspace{0.1cm}
		\textbf{DETAILED TUTORIAL:}
		\begin{itemize}
			\item \textbf{Context-dependent stop words:}
			\begin{itemize}
				\item Some words may be useful in specific domains
				\item "bank" is common but important in financial text
				\item Standard lists may remove important words
			\end{itemize}
		\end{itemize}
		
		\textcolor{myred}{\textbf{EXAM TRAP:}} Stop words are \textcolor{myblue}{context-dependent} - not one-size-fits-all!
	\end{frame}
	
	\begin{frame}{Q94: DTM Example - Tutorial}
		\fontsize{6.5}{7.5}\selectfont
		\textbf{QUESTION: In Figure 9.1, what does the Document Term Matrix show?}
		
		\vspace{0.1cm}
		\textbf{OPTIONS:}
		\begin{enumerate}
			\item Document titles and authors
			\item Rows as documents, columns as words, entries as word counts
			\item Document dates
			\item Document sources
		\end{enumerate}
		
		\vspace{0.1cm}
		\textcolor{myblue}{\textbf{ANSWER: B}}
		
		\vspace{0.1cm}
		\textbf{DETAILED TUTORIAL:}
		\begin{itemize}
			\item \textbf{DTM structure:}
			\begin{itemize}
				\item \textbf{Rows:} Each row = one document
				\item \textbf{Columns:} Each column = one word from vocabulary
				\item \textbf{Entries:} Word counts or frequencies
			\end{itemize}
		\end{itemize}
		
		\textcolor{myred}{\textbf{EXAM TRAP:}} DTM = \textcolor{myblue}{documents $\times$ words matrix with counts}!
	\end{frame}
	
	\begin{frame}{Q95: Feature Extraction Purpose - Tutorial}
		\fontsize{6.5}{7.5}\selectfont
		\textbf{QUESTION: Why is feature extraction (text representation) important in NLP?}
		
		\vspace{0.1cm}
		\textbf{OPTIONS:}
		\begin{enumerate}
			\item To make text longer
			\item To convert text into numeric vectors for machine learning analysis
			\item To remove words
			\item To make text more readable
		\end{enumerate}
		
		\vspace{0.1cm}
		\textcolor{myblue}{\textbf{ANSWER: B}}
		
		\vspace{0.1cm}
		\textbf{DETAILED TUTORIAL:}
		\begin{itemize}
			\item \textbf{Feature extraction purpose:}
			\begin{itemize}
				\item Machine learning models need numeric inputs
				\item Text must be converted to numbers
				\item Enables statistical analysis
			\end{itemize}
		\end{itemize}
		
		\textcolor{myred}{\textbf{EXAM TRAP:}} Feature extraction = \textcolor{myblue}{text $\to$ numeric vectors}!
	\end{frame}
	
	\begin{frame}{Q96: Document Term Matrix Dimensions - Tutorial}
		\fontsize{6.5}{7.5}\selectfont
		\textbf{QUESTION: What are the dimensions of a Document Term Matrix?}
		
		\vspace{0.1cm}
		\textbf{OPTIONS:}
		\begin{enumerate}
			\item $D \times W$ (documents $\times$ words)
			\item $W \times D$ (words $\times$ documents)
			\item $D \times D$ (documents $\times$ documents)
			\item $W \times W$ (words $\times$ words)
		\end{enumerate}
		
		\vspace{0.1cm}
		\textcolor{myblue}{\textbf{ANSWER: A}}
		
		\vspace{0.1cm}
		\textbf{DETAILED TUTORIAL:}
		\begin{itemize}
			\item \textbf{DTM dimensions:}
			\begin{itemize}
				\item Rows: D = number of documents
				\item Columns: W = number of words in vocabulary
				\item Dimensions: D $\times$ W
			\end{itemize}
		\end{itemize}
		
		\textcolor{myred}{\textbf{EXAM TRAP:}} DTM dimensions = \textcolor{myblue}{documents $\times$ words}!
	\end{frame}
	
	\begin{frame}{Q97: BoW Vector Example - Tutorial}
		\fontsize{6.5}{7.5}\selectfont
		\textbf{QUESTION: In the BoW example, why does d1 have a 2 for "high" in the count vector?}
		
		\vspace{0.1cm}
		\textbf{OPTIONS:}
		\begin{enumerate}
			\item Because "high" is the most important word
			\item Because "high" appears twice in the document
			\item Because "high" is a stop word
			\item Because "high" is the first word
		\end{enumerate}
		
		\vspace{0.1cm}
		\textcolor{myblue}{\textbf{ANSWER: B}}
		
		\vspace{0.1cm}
		\textbf{DETAILED TUTORIAL:}
		\begin{itemize}
			\item \textbf{Example text:}
			\begin{itemize}
				\item d1: "... reaching a high of 1550 ... with a high of 1560"
				\item "high" appears twice
				\item Count vector records 2
			\end{itemize}
		\end{itemize}
		
		\textcolor{myred}{\textbf{EXAM TRAP:}} Count BoW = \textcolor{myblue}{actual frequency} - not just presence!
	\end{frame}
	
	\begin{frame}{Q98: Bag of n-grams - Tutorial}
		\fontsize{6.5}{7.5}\selectfont
		\textbf{QUESTION: What is the Bag of n-grams (BoN) approach?}
		
		\vspace{0.1cm}
		\textbf{OPTIONS:}
		\begin{enumerate}
			\item A method to count single words
			\item Applying BoW to n-grams (n contiguous words as units)
			\item A method to remove n-grams
			\item A method to count only bigrams
		\end{enumerate}
		
		\vspace{0.1cm}
		\textcolor{myblue}{\textbf{ANSWER: B}}
		
		\vspace{0.1cm}
		\textbf{DETAILED TUTORIAL:}
		\begin{itemize}
			\item \textbf{BoN definition:}
			\begin{itemize}
				\item Extends BoW to n-grams
				\item n contiguous words = one unit
				\item Handles phrases and negations
			\end{itemize}
		\end{itemize}
		
		\textcolor{myred}{\textbf{EXAM TRAP:}} BoN = \textcolor{myblue}{BoW applied to n-grams} - captures phrases!
	\end{frame}
	
	\begin{frame}{Q99: Cosine Similarity Application - Tutorial}
		\fontsize{6.5}{7.5}\selectfont
		\textbf{QUESTION: In document similarity, what does a cosine similarity close to 1 indicate?}
		
		\vspace{0.1cm}
		\textbf{OPTIONS:}
		\begin{enumerate}
			\item Documents are very different
			\item Documents are very similar (vectors point in same direction)
			\item Documents are opposite
			\item Documents are unrelated
		\end{enumerate}
		
		\vspace{0.1cm}
		\textcolor{myblue}{\textbf{ANSWER: B}}
		
		\vspace{0.1cm}
		\textbf{DETAILED TUTORIAL:}
		\begin{itemize}
			\item \textbf{Cosine similarity interpretation:}
			\begin{itemize}
				\item cos equals 1: Very similar documents
				\item cos equals 0: Unrelated documents
				\item cos equals -1: Opposite documents
			\end{itemize}
		\end{itemize}
		
		\textcolor{myred}{\textbf{EXAM TRAP:}} cos equals 1 = \textcolor{myblue}{very similar documents}!
	\end{frame}
	
	\begin{frame}{Q100: NLP Chapter Summary - Tutorial}
		\fontsize{6.5}{7.5}\selectfont
		\textbf{QUESTION: What are the key takeaways from the NLP chapter?}
		
		\vspace{0.1cm}
		\textbf{OPTIONS:}
		\begin{enumerate}
			\item Only BoW is used in NLP
			\item NLP involves pre-processing, feature extraction, dictionary or ML approaches, and careful evaluation
			\item Only machine learning is used
			\item NLP doesn't require pre-processing
		\end{enumerate}
		
		\vspace{0.1cm}
		\textcolor{myblue}{\textbf{ANSWER: B}}
		
		\vspace{0.1cm}
		\textbf{DETAILED TUTORIAL:}
		\begin{itemize}
			\item \textbf{Key takeaways:}
			\begin{enumerate}
				\item \textbf{Pre-processing:} Tokenization, stop words, stemming/lemmatization
				\item \textbf{Feature extraction:} BoW, n-grams, TF-IDF
				\item \textbf{Approaches:} Dictionary-based and machine learning
				\item \textbf{Evaluation:} Confusion matrix metrics, speed, data requirements
			\end{enumerate}
		\end{itemize}
		
		\textcolor{myred}{\textbf{EXAM TRAP:}} NLP = \textcolor{myblue}{pre-processing + feature extraction + modeling + evaluation}!
	\end{frame}
	
	% ============================================================
	% SUMMARY SLIDE
	% ============================================================
	
	\begin{frame}{Summary: Complete NLP Review}
		\fontsize{6.5}{7.5}\selectfont
		\textbf{\textcolor{myblue}{Pre-processing:}}
		\begin{itemize}
			\item \textcolor{myblue}{Tokenization:} Sentences $\to$ words
			\item \textcolor{myblue}{Stop Words:} Remove common words with no value
			\item \textcolor{myblue}{Stemming:} Cut prefixes/suffixes (fast, less accurate)
			\item \textcolor{myblue}{Lemmatization:} Dictionary form (slow, more accurate)
		\end{itemize}
		
		\vspace{0.1cm}
		\textbf{\textcolor{myblue}{Feature Extraction:}}
		\begin{itemize}
			\item \textcolor{myblue}{BoW:} Binary (presence) or Count (frequency)
			\item \textcolor{myblue}{n-grams:} Contiguous word sequences
			\item \textcolor{myblue}{TF-IDF:} $TF \times IDF$ - weights rare words higher
			\item \textcolor{myblue}{Embeddings:} Dense vectors capturing meaning (Word2Vec)
		\end{itemize}
		
		\vspace{0.1cm}
		\textbf{\textcolor{myblue}{Classification:}}
		\begin{itemize}
			\item \textcolor{myblue}{Dictionary:} Pre-defined word lists (sentiment analysis)
			\item \textcolor{myblue}{Naïve Bayes:} Simple, fast, generative classifier
			\item \textcolor{myblue}{Smoothing:} Laplace smoothing to avoid zero probabilities
		\end{itemize}
		
		\vspace{0.1cm}
		\textbf{\textcolor{myblue}{Evaluation:}}
		\begin{itemize}
			\item \textcolor{myblue}{Confusion Matrix:} TP, TN, FP, FN
			\item \textcolor{myblue}{Metrics:} Accuracy, Precision, Recall, F1
			\item \textcolor{myblue}{Challenges:} Ambiguity, context, negation, sarcasm
		\end{itemize}
		
		\textcolor{myred}{\textbf{Exam Success:}} Understand the process from raw text to classification and evaluation!
	\end{frame}
	
\end{document}